\section{Orders of degeneration at the nef boundary}
\label{sec:nef-volume-degeneration}

If $\theta$ is nef and $h$ is Hermitian, the volumes $\mathcal V_\mu(\theta+\epsilon h)$ are positive and, by \cref{thm:gaud-volume-nef-limit}, decrease to $\mathcal V_\mu(\theta)$ as $\epsilon\downarrow0$. When this limit is zero, the rate of decay measures how degenerate $\theta$ is at the nef boundary. For a closed nef class on a compact K\"ahler manifold, the volume of the perturbation is the intersection polynomial $\int_X(\theta+\epsilon h)^n$, so it decays like an integral power of $\epsilon$ determined by the numerical dimension. For a nonclosed nef background the logarithmic rate of decay can be fractional, and it need not converge. We therefore define lower and upper orders of degeneration as lower and upper limits of this rate. Changing the probability measure or the perturbing metric only introduces bounded factors, so the orders are invariants of the background.

Let $X$ again be a connected compact complex manifold of dimension $n\geq1$, and $\mu$ a smooth strictly positive probability form on $X$. Let $\theta$ be nef and let $h$ be Hermitian. For $0<\epsilon<1$, put
\[
 a_{\mu,h}(\theta;\epsilon)
   =\frac{\log\mathcal V_\mu(\theta+\epsilon h)}{\log\epsilon}.
\]
The logarithm is finite because $\theta+\epsilon h$ belongs to $\mathscr P$, so that its volume is positive.

\Needspace{16\baselineskip}
\begin{proposition}\label{prop:nef-degeneration-orders}\index{Nef degeneration!orders}
The limits
\[
 \underline a(\theta)=\liminf_{\epsilon\downarrow0}
                         a_{\mu,h}(\theta;\epsilon),
 \qquad
 \overline a(\theta)=\limsup_{\epsilon\downarrow0}
                         a_{\mu,h}(\theta;\epsilon)
\]
are independent of $\mu$ and $h$, and satisfy
\[
 0\leq\underline a(\theta)\leq\overline a(\theta)\leq n.
\]
They are unchanged when $\theta$ is replaced by $\theta+\ddc f$, with $f$ smooth, or by $r\theta$, with $r>0$. If $\mathcal V_\mu(\theta)>0$, both limits are zero.
\end{proposition}
\begin{proof}
Since $\theta$ is pseudoeffective, monotonicity and homogeneity give
\[
 \epsilon\lambda_\mu(h)
       \leq\lambda_\mu(\theta+\epsilon h)
       \leq\lambda_\mu(\theta+h),
 \qquad 0<\epsilon\leq1.
\]
Thus there are constants $c,C>0$ such that
\[
 c\epsilon^n\leq\mathcal V_\mu(\theta+\epsilon h)\leq C,
 \qquad 0<\epsilon\leq1.
\]
Taking logarithms and dividing by the negative number $\log\epsilon$ shows that the quotients $a_{\mu,h}(\theta;\epsilon)$ stay bounded as $\epsilon\downarrow0$, and that their lower and upper limits lie in $[0,n]$.

By \cref{prop:gaud-volume-density-comparison}, changing the smooth probability measure multiplies the volume by a factor between two fixed positive constants. The resulting error in the logarithm is bounded, and after division by $\log\epsilon$ it tends to zero.

\Needspace{8\baselineskip}
For another Hermitian metric $h'$, choose constants $c_2\geq c_1>0$ with $c_1h\leq h'\leq c_2h$. Then
\[
 \mathcal V_\mu(\theta+c_1\epsilon h)
 \leq\mathcal V_\mu(\theta+\epsilon h')
 \leq\mathcal V_\mu(\theta+c_2\epsilon h).
\]
For any fixed $t>0$ and small $\epsilon$,
\[
 \frac{\log\mathcal V_\mu(\theta+t\epsilon h)}{\log\epsilon}
 =a_{\mu,h}(\theta;t\epsilon)
       \frac{\log(t\epsilon)}{\log\epsilon}.
\]
The last factor tends to one and the first is bounded, while the rescaling $\epsilon\mapsto t\epsilon$ does not change lower and upper limits as $\epsilon\downarrow0$. Hence the left side has the same lower and upper limits as $a_{\mu,h}(\theta;\epsilon)$. Dividing the preceding inequalities by $\log\epsilon<0$ reverses them, and taking lower and upper limits proves independence of the metric.

Invariance under smooth $\ddc$ changes follows from that of $\mathcal V_\mu$. For scaling, homogeneity gives
\[
 \mathcal V_\mu(r\theta+\epsilon h)
       =r^n\mathcal V_\mu(\theta+(\epsilon/r)h).
\]
The bounded error $n\log r$ and the same rescaling argument prove scale invariance. Finally, if $\mathcal V_\mu(\theta)>0$, then by \cref{thm:gaud-volume-nef-limit} the perturbed volumes lie between $\mathcal V_\mu(\theta)$ and $\mathcal V_\mu(\theta+h)$ for $0<\epsilon\leq1$, so both limits are zero.
\end{proof}

We call $\underline a(\theta)$ and $\overline a(\theta)$ the \emph{lower} and \emph{upper orders of degeneration} of the nef background $\theta$. When they agree, their common value is the \emph{common order}, and we then also say that $\theta$ has this order. The two orders can differ, and a common order can be irrational. Even when both are zero, the volume can tend to zero, at a logarithmic rate. The elliptic examples in \cref{sec:nef-surface-orders} exhibit each of these behaviors. Closed nef classes on K\"ahler manifolds, in contrast, degenerate polynomially.

\Needspace{15\baselineskip}
\begin{proposition}\label{prop:nef-degeneration-kahler}
Suppose $X$ is K\"ahler, $\theta$ is closed and nef, and $h$ is K\"ahler. Set
\[
 r=\max\left\{0\leq k\leq n:
          \int_X\theta^k\wedge h^{n-k}>0\right\}.
\]
Then
\[
 \underline a(\theta)=\overline a(\theta)=n-r.
\]
Thus $n$ minus the common order is the numerical dimension of the closed nef class of $\theta$.
\end{proposition}
\begin{proof}
All the intersection numbers in the definition of $r$ are nonnegative. Indeed, the mixed integrals against $h$ of smooth positive representatives of $\theta+\delta h$ are positive, agree with the cohomological integrals by closed Stokes, and converge to the stated numbers as $\delta\downarrow0$. The number for $k=0$ is $\int_Xh^n>0$, so $r$ is well defined.

By the closed nef formula following \cref{thm:gaud-volume-nef-limit},
\[
 \mathcal V_\mu(\theta+\epsilon h)
       =\int_X(\theta+\epsilon h)^n
       =\sum_{k=0}^n\binom nk\epsilon^{n-k}
                         \int_X\theta^k\wedge h^{n-k}.
\]
As $\epsilon\downarrow0$, the leading nonzero term is a positive constant times $\epsilon^{n-r}$, and taking the logarithmic quotient proves the assertion. By \cref{prop:nef-degeneration-orders}, the result holds for every Hermitian perturbation, not only for K\"ahler ones.
\end{proof}

\Needspace{17\baselineskip}
The orders can also be computed on the non-K\"ahler Hopf threefold $X$ of \cref{chap:hopf}. Keep its forms $\omega=\alpha+\beta$, where $\alpha\geq0$ is closed of rank two, $\beta\geq0$ has rank one, $\alpha^3=\beta^2=0$, and $\int_X\omega^3=6$. Take $\theta=\alpha$, $h=\omega$, and $\mu_0=\omega^3/6$. The metric
\[
 \Omega_\epsilon=\alpha+\epsilon\omega
                    =(1+\epsilon)\alpha+\epsilon\beta
\]
is strictly positive. Its volume is
\[
 \Omega_\epsilon^3=\epsilon(1+\epsilon)^2\omega^3
                  =6\epsilon(1+\epsilon)^2\mu_0.
\]
Hence the zero potential solves the smooth normalized equation, and
\[
 \mathcal V_{\mu_0}(\alpha+\epsilon\omega)
             =6\epsilon(1+\epsilon)^2,
 \qquad
 \underline a(\alpha)=\overline a(\alpha)=1.
\]
By \cref{prop:nef-degeneration-orders}, the same orders are obtained for every smooth positive probability measure and every Hermitian perturbation.

Each perturbed metric nevertheless has infinite upper mass. If $V_j$ are the smooth potentials of \cref{thm:hopf-unbounded}, then
\[
 \Omega_\epsilon+\ddc(\epsilon V_j)
       =\alpha+\epsilon(\omega+\ddc V_j)>0,
\]
and monotonicity of the determinant of positive matrices gives
\[
 \int_X\bigl(\Omega_\epsilon+\ddc(\epsilon V_j)\bigr)^3
       \geq\epsilon^3\int_X(\omega+\ddc V_j)^3\longrightarrow+\infty.
\]
Thus a finite order of degeneration at the boundary coexists with infinite upper mass at every positive perturbation.

\subsection{Elliptic surface volumes and order spectra}\label{sec:nef-surface-orders}

A coefficient that vanishes along one elliptic factor determines the harmonic mean in the volume formula. Finite-order zeros give fractional or logarithmic decay; flat coefficients allow the logarithmic slope to oscillate.

Let $E=\mathbb C/(\mathbb Z+\mathrm{i}\mathbb Z)$, write $X=E_z\times E_w$, and fix
\[
 \omega_1=\frac{\mathrm{i}}{2}\mathrm{d} z\wedge\mathrm{d}\bar z,\qquad
 \omega_2=\frac{\mathrm{i}}{2}\mathrm{d} w\wedge\mathrm{d}\bar w,\qquad
 \mu=\omega_1\wedge\omega_2,\qquad h=\omega_1+\omega_2.
\]
Both area forms have integral one. Let $a,q\geq0$ and $r$ be smooth functions of $z$ alone such that $|r|^2\leq aq$, and set
\[
 \theta=\frac{\mathrm{i}}{2}\bigl(a\,\mathrm{d} z\wedge\mathrm{d}\bar z+
 r\,\mathrm{d} z\wedge\mathrm{d}\bar w+\bar r\,\mathrm{d} w\wedge\mathrm{d}\bar z+
 q\,\mathrm{d} w\wedge\mathrm{d}\bar w\bigr).
\]
This is a smooth semipositive, hence nef, form.

\begin{proposition}\label{prop:nef-elliptic-exact-volume}
For $\epsilon>0$, put $g=q+\epsilon$ and $I=\int_Eg^{-1}\omega_1$. Then
\begin{equation}\label{eq:nef-elliptic-volume}
 \mathcal V_\mu(\theta+\epsilon h)
 =\frac{2}{I}\left(\int_E(a+\epsilon)\omega_1-
 \int_E\frac{|r|^2}{g}\omega_1\right).
\end{equation}
The infimum defining this volume is attained by a smooth positive Gauduchon metric. In particular, if $a=b>0$ and $r=0$, then $\mathcal V_\mu(\theta+\epsilon h)=2(b+\epsilon)/I$. If $|r|^2=aq$, then
\begin{equation}\label{eq:nef-rank-one-elliptic-volume}
 \mathcal V_\mu(\theta+\epsilon h)=\frac{2\epsilon(1+K)}{I},
 \qquad K=\int_E\frac{a}{q+\epsilon}\omega_1.
\end{equation}
\end{proposition}

\begin{proof}
Define
\[
 c=I^{-1}\left(\int_E(a+\epsilon)\omega_1-
 \int_E\frac{|r|^2}{g}\omega_1\right),\qquad
 A=\frac{c+|r|^2}{g}.
\]
The inequality $(a+\epsilon)g-|r|^2>0$ gives $c>0$, while the definition of $c$ makes $A-a-\epsilon$ have zero $\omega_1$-mean. Division of its nonconstant Fourier coefficients by the nonzero flat Laplace eigenvalues produces a smooth real function $u$ on $E$ with $\ddc u=(A-a-\epsilon)\omega_1$. Pull it back to $X$. The form $\Omega=\theta+\epsilon h+\ddc u$ has coefficient matrix $\left(\begin{smallmatrix}A&r\\\bar r&g\end{smallmatrix}\right)$, determinant $c$, and $\Omega^2=2c\mu$.

The metric $\gamma_*=\Omega/g$ has constant fiber coefficient $1$; its remaining coefficients depend only on $z$. Applying $\ddc$ to its base or mixed terms repeats $\mathrm{d} z$ or $\mathrm{d}\bar z$, so $\ddc\gamma_*=0$. It is smooth, positive, and pointwise conformal to $\Omega$. The arithmetic--geometric mean argument in \cref{thm:gaud-volume-dual-formula} and equality for $\gamma_*$ prove \eqref{eq:nef-elliptic-volume}. The two specializations follow by direct substitution.
\end{proof}

\begin{proposition}\label{prop:nef-diagonal-fractional-and-zero}\index{Nef degeneration!fractional and zero orders}
There are smooth semipositive nonclosed forms on the fixed surface $X$ with common order $1-1/m$ for each integer $m>1$. There is also one with
\[
 \mathcal V_\mu(\theta+\epsilon h)\sim\frac{2b}{\pi\log(1/\epsilon)}
 \quad(\epsilon\downarrow0),
\]
for any chosen $b>0$; it has boundary volume zero and both orders zero.
\end{proposition}

\begin{proof}
In \cref{prop:nef-elliptic-exact-volume} take $a=b>0$ and $r=0$. Choose $q$ smooth, positive outside $z_0\in E$, and equal to $|\zeta|^{2m}$ in a small flat coordinate disk about $z_0$. A positive smooth extension exists away from the disk. For $m=1$, polar coordinates give
\[
 I=\int_E\frac{\omega_1}{q+\epsilon}
   =\pi\log(1/\epsilon)+O(1).
\]
Equation \eqref{eq:nef-elliptic-volume} gives the displayed asymptotic. If $m>1$, the change $s=|\zeta|^2=\epsilon^{1/m}t$ gives
\[
 I\sim C_m\epsilon^{1/m-1},\qquad
 C_m=\pi\int_0^\infty\frac{\mathrm{d} t}{1+t^m}>0,
\]
and the common order is $1-1/m$. In the logarithmic case $\lambda_\mu(\theta)\geq0$ by semipositivity, while $\lambda_\mu(\theta)\leq\lambda_\mu(\theta+\epsilon h)\to0$ by monotonicity. Hence $\mathcal V_\mu(\theta)=0$. The quotient of the displayed volume by $\epsilon^\delta$ diverges for every $\delta>0$. Since $\mathrm{d}\theta=\mathrm{d} q\wedge\omega_2\ne0$, these backgrounds are nonclosed.
\end{proof}

\begin{proposition}\label{prop:nef-oscillatory-orders}\index{Nef degeneration!unequal orders}
For every $0<\alpha<\beta<1$ there is a smooth semipositive nonclosed background on the same $(X,\mu,h)$ with
\[
 \bigl(\underline a(\theta),\overline a(\theta)\bigr)
 =(1+\alpha,1+\beta).
\]
For every $0<\sigma<1$ there is such a background with common order $1+\sigma$.
\end{proposition}

\begin{proof}
For all sufficiently large $t$, define
\[
 p(t)=\frac{\alpha+\beta}{2}
 +\frac{\beta-\alpha}{2}\sin(\log\log t),\qquad P(t)=tp(t).
\]
Choose $T$ large enough that $0<c_0\leq P'(t)\leq d_0<1$ for $t\geq T$; this follows from
\[
 P'(t)=p(t)+\frac{\beta-\alpha}{2\log t}\cos(\log\log t).
\]
On $E\setminus\{z_0\}$ choose a smooth $t(z)\geq T$ equal to $|\zeta|^{-2}$ near $z_0$ and bounded outside every neighborhood of $z_0$. Put
\[
 q=\mathrm{e}^{-t},\qquad a=\mathrm{e}^{-P(t)},\qquad
 r=\mathrm{e}^{-(P(t)+t)/2},
\]
and set all three equal to zero at $z_0$. These functions are smooth and flat there: $P(t)\geq\alpha t$ and every fixed derivative of $P$ grows at most polynomially in $t$, so the chain rule bounds each coordinate derivative by a polynomial in $|\zeta|^{-1}$ times $\mathrm{e}^{-c/|\zeta|^2}$. We have $|r|^2=aq$, so the resulting $\theta$ is semipositive. It is nonclosed: the component $\mathrm{d} q\wedge\omega_2$ of $\mathrm{d}\theta$ has two fiber differentials, which derivatives of the remaining terms cannot cancel.

Write $\epsilon=\mathrm{e}^{-L}$. Choose a fixed integration endpoint $T_0\geq T$ inside the coordinate disk, with the transition to the bounded global $t(z)$ outside a smaller disk. The complementary contributions to $I$ and $K$ in \eqref{eq:nef-rank-one-elliptic-volume} are $O(1)$, while polar coordinates on the disk give
\[
 I_{\mathrm{loc}}=\pi\int_{T_0}^\infty
 \frac{t^{-2}\,\mathrm{d} t}{\mathrm{e}^{-t}+\mathrm{e}^{-L}},\qquad
 K_{\mathrm{loc}}=\pi\int_{T_0}^\infty
 \frac{\mathrm{e}^{-P(t)}t^{-2}\,\mathrm{d} t}{\mathrm{e}^{-t}+\mathrm{e}^{-L}}.
\]
Integrating $I$ over $t\geq L$ gives a lower bound by a fixed multiple of $\mathrm{e}^L/L$. For the upper bound, use $(\mathrm{e}^{-t}+\mathrm{e}^{-L})^{-1}\leq\mathrm{e}^t$ for $t\leq L$ and at most $\mathrm{e}^L$ for $t\geq L$, then split $[T_0,L]$ at $L/2$. Thus
\begin{equation}\label{eq:nef-elliptic-I-asymptotic}
 \log I=L-\log L+O(1).
\end{equation}
The derivatives of $P$ and $t-P(t)$ are positive and bounded away from zero. Splitting the $K$-integral at $L$ therefore gives $K\leq C\mathrm{e}^{L-P(L)}$. On $[L,L+1]$ its denominator is at most $2\mathrm{e}^{-L}$ and $P(t)\leq P(L)+1$, so
\[
 c_1L^{-2}\mathrm{e}^{L-P(L)}\leq K\leq c_2\mathrm{e}^{L-P(L)}.
\]
The complementary $O(1)$ term is absorbed because $P(L)\leq\beta L$. Consequently $\log K=L-P(L)+O(\log L)$ and $K\to\infty$. Combining this with \eqref{eq:nef-rank-one-elliptic-volume} and \eqref{eq:nef-elliptic-I-asymptotic} yields
\begin{equation}\label{eq:nef-oscillatory-log-volume}
 \log \mathcal V_\mu(\theta+\mathrm{e}^{-L}h)=-L-P(L)+O(\log L),
 \qquad
 \frac{\log \mathcal V_\mu(\theta+\epsilon h)}{\log\epsilon}
 =1+p(L)+o(1).
\end{equation}
The sine in $p(L)$ has limits $-1$ and $1$ along sequences, proving the first assertion. If instead $P(t)=\sigma t$, the same smoothness and integral estimates give common order $1+\sigma$. In particular $\sigma=\sqrt2-1$ gives an irrational common order.
\end{proof}

\subsection{Comparison under an arbitrary modification}\label{sec:nef-modification-comparison}

The orders of degeneration are defined on a fixed manifold. We now show that pulling a nef background back by a modification can only increase them. Recall that a modification $\pi:Y\to X$ is a proper surjective holomorphic map that restricts to a biholomorphism between dense Zariski open subsets. The proof compares Gauduchon tests on $X$ with tests on $Y$. The pullback of a Gauduchon test is degenerate along the exceptional locus, and adding a small positive test on $Y$ repairs this. To control the effect of the added test uniformly, we first bound the pairing of pulled-back tests with a Hermitian form on $Y$ by their pairing with a Hermitian form on $X$, uniformly over the whole cone of tests.

\Needspace{8\baselineskip}
In this subsection $\pi:Y\to X$ is a modification of connected compact complex manifolds of dimension $n\geq2$. We use Hironaka's analytic elimination of indeterminacies in the following form: there are a finite composition $p:Z\to X$ of blowups along smooth centers and a modification $F:Z\to Y$ with $p=\pi\circ F$. It is obtained by resolving the indeterminacies of the inverse meromorphic map $\pi^{-1}\circ p$; the identity $p=\pi\circ F$ holds on the dense set where the maps are biholomorphic, hence everywhere. This classical input is discussed in \cref{sec:history-gaud-volume}. For an effective real divisor $D$, we write $\eta_D$ for the curvature form of a smooth Hermitian metric on the corresponding real line bundle, as in \cref{sec:divisorial-gaud-volume}.

\begin{lemma}\label{lem:nef-blowup-metric}
Let $p:Z\to X$ be a finite composition of blowups along smooth centers of codimension at least two, and let $h_X$ be a Hermitian form on $X$. There are a constant $C_0>0$, an effective real divisor $D$ on $Z$ whose components are exceptional for $p$, and a smooth curvature form $\eta_D$ such that
\begin{equation}\label{eq:nef-blowup-positive-metric}
 h_Z=C_0p^*h_X-\eta_D>0.
\end{equation}
Every component of $D$ is mapped by $p$ into a set of complex codimension at least two in $X$.
\end{lemma}

\begin{proof}
First let $q:M'\to M$ be one blowup along a smooth center, with exceptional divisor $E$, and let $h_M$ be a Hermitian form on $M$. The restriction of $\mathcal O(-E)$ to $E$ is the tautological line bundle of the projectivized normal bundle of the center, which is positive on the projective fibers. A Hermitian metric on that normal bundle therefore induces a metric on $\mathcal O(-E)|_E$ whose curvature is positive on the tangent spaces of the fibers. Extend this metric to $\mathcal O(-E)$ on $M'$ by extending the smooth logarithm of its ratio to a fixed reference metric from the closed submanifold $E$, and let $\beta$ be its curvature. The semipositive form $q^*h_M$ is positive off $E$, and at points of $E$ its kernel is exactly the tangent space of the fiber, on which $\beta$ is positive. Hence $Cq^*h_M+\beta>0$ for every sufficiently large $C$: otherwise, unit tangent vectors violating the inequality for $C\to\infty$ would converge, by compactness, to a vector in the kernel of $q^*h_M$ on which $\beta$ is nonpositive. With $\eta_E=-\beta$, the curvature of the induced metric on $\mathcal O(E)$, this reads $Cq^*h_M-\eta_E>0$.

For a composition of blowups, apply this step successively, with $h_M$ the positive form already obtained on the previous blowup. Each step multiplies the previous divisor by a positive constant, pulls it back and adds the new exceptional divisor, so the divisor stays effective and exceptional. This proves \eqref{eq:nef-blowup-positive-metric}. Each exceptional divisor maps into the image of its center, which has dimension at most $n-2$.
\end{proof}

\begin{lemma}\label{lem:nef-uniform-pairing}
Fix Hermitian forms $h_X,h_Y$ for a modification $\pi:Y\to X$. There is a constant $K>0$ such that
\begin{equation}\label{eq:nef-uniform-pairing}
 \int_Y h_Y\wedge\pi^*\Psi\leq K\int_X h_X\wedge\Psi
\end{equation}
for every smooth strictly positive $(n-1,n-1)$-form $\Psi$ on $X$ with $\ddc\Psi=0$.
\end{lemma}

\begin{proof}
Use the factorization $p=\pi\circ F$ above and \cref{lem:nef-blowup-metric}. By compactness there is $a>0$ with $h_Z\geq aF^*h_Y$. Pulling back the semipositive form $\Psi$ and integrating, we obtain
\[
 a\int_Yh_Y\wedge\pi^*\Psi
 \leq\int_Zh_Z\wedge p^*\Psi
 =C_0\int_Xh_X\wedge\Psi-\int_Dp^*\Psi.
\]
Here the changes of variables use that $F$ and $p$ have degree one, and the Poincar\'e--Lelong formula, together with $\ddc p^*\Psi=0$, replaces the pairing with $\eta_D$ by the integral over $D$. This last integral vanishes: on the regular locus of each component of $D$, the restriction of $p$ has rank at most $n-2$, so the restriction of the $(n-1,n-1)$-form $p^*\Psi$ is zero. Taking $K=C_0/a$ proves the assertion.
\end{proof}

\Needspace{17\baselineskip}
\begin{proposition}\label{prop:nef-modification-volume}\index{Nef degeneration!modification}
Let $\pi:Y\to X$ be a modification of connected compact complex manifolds of dimension $n\geq2$, and let $\theta$ be a smooth nef background on $X$. Then $\pi^*\theta$ is nef. For fixed Hermitian forms $h_X,h_Y$ and smooth positive probability forms $\mu_X,\mu_Y$, there are constants $C,K>0$, independent of $\epsilon$, such that
\begin{equation}\label{eq:nef-main-modification-comparison}
 \mathcal V_{\mu_Y}(\pi^*\theta+\epsilon h_Y)
 \leq C\mathcal V_{\mu_X}(\theta+K\epsilon h_X)
 \qquad(\epsilon>0).
\end{equation}
Consequently
\[
 \underline a(\pi^*\theta)\geq\underline a(\theta),\qquad
 \overline a(\pi^*\theta)\geq\overline a(\theta).
\]
In dimension one a modification is an isomorphism and preserves both orders.
\end{proposition}

\begin{proof}
\emph{Step 1: nefness.} Choose $C'$ with $\pi^*h_X\leq C'h_Y$. For $\epsilon>0$, pull back a smooth representative $\theta+\ddc u_\tau\geq-\tau h_X$ with $\tau<\epsilon/C'$; this gives a positive representative of $\pi^*\theta+\epsilon h_Y$. Thus $\pi^*\theta$ is nef, and both volumes in \eqref{eq:nef-main-modification-comparison} are positive.

\emph{Step 2: regularized tests.} The ratio $\pi^*\mu_X/\mu_Y$ is smooth, nonnegative and bounded on the compact manifold $Y$, so there is $c>0$ with $\mu_Y\geq c\pi^*\mu_X$. Fix a smooth Gauduchon metric $\gamma_0$ on $Y$, given by \cref{thm:gauduchon}. For an arbitrary Gauduchon metric $\gamma$ on $X$, write $\Psi=\gamma^{n-1}$, and for $\delta>0$ define $\Gamma_\delta$ by
\[
 \Gamma_\delta^{n-1}=\pi^*\Psi+\delta\gamma_0^{n-1}.
\]
The right side is a strictly positive $(n-1,n-1)$-form, and such a form has a unique smooth positive $(n-1)$st root, obtained by solving for the Hermitian matrix of the root from the positive cofactor matrix. Since the right side is $\ddc$-closed, $\Gamma_\delta$ is a smooth Gauduchon metric on $Y$.

\emph{Step 3: the numerator.} For $\epsilon>0$, \cref{lem:nef-uniform-pairing} gives
\begin{equation}\label{eq:nef-modification-numerator}
 \lim_{\delta\downarrow0}\int_Y(\pi^*\theta+\epsilon h_Y)
 \wedge\Gamma_\delta^{n-1}
 =\int_X\theta\wedge\Psi+\epsilon\int_Yh_Y\wedge\pi^*\Psi
 \leq\int_X(\theta+K\epsilon h_X)\wedge\Psi.
\end{equation}
The first term in the middle is nonnegative because $\theta$ is nef, and the second is positive because $\pi^*\Psi$ is positive on the set where $\pi$ is biholomorphic, so the limit is positive.

\emph{Step 4: the denominator.} We claim that
\begin{equation}\label{eq:nef-modification-denominator}
 \lim_{\delta\downarrow0}D_{\mu_Y}(\Gamma_\delta)
 =\int_Y(\pi^*\gamma^n)^{(n-1)/n}\mu_Y^{1/n}
 \geq c^{1/n}D_{\mu_X}(\gamma)>0.
\end{equation}
To justify the limit on the exceptional locus, represent $\Gamma_\delta^{n-1}$ by its positive cofactor matrix. The density of $(\Gamma_\delta^n)^{(n-1)/n}$ is a fixed dimensional constant times the $1/n$th power of the determinant of that matrix, and this expression extends continuously to positive semidefinite matrices. Since $Y$ is compact, dominated convergence applies. The inequality follows from $\mu_Y\geq c\pi^*\mu_X$ and the change of variables for the degree-one map $\pi$.

\emph{Step 5: conclusion.} Test the infimum defining $\lambda_{\mu_Y}(\pi^*\theta+\epsilon h_Y)$ with $\Gamma_\delta$, use \eqref{eq:nef-modification-numerator} and \eqref{eq:nef-modification-denominator}, and let $\delta\downarrow0$. For every Gauduchon metric $\gamma$ on $X$ this gives
\[
 \lambda_{\mu_Y}(\pi^*\theta+\epsilon h_Y)
 \leq c^{-1/n}
 \frac{\int_X(\theta+K\epsilon h_X)\wedge\gamma^{n-1}}{D_{\mu_X}(\gamma)}.
\]
Taking the infimum over $\gamma$ and then the $n$th power proves \eqref{eq:nef-main-modification-comparison} with $C=c^{-1}$.

Dividing the logarithm of \eqref{eq:nef-main-modification-comparison} by $\log\epsilon<0$ reverses the inequality. The bounded error $\log C$ and the rescaling $\epsilon\mapsto K\epsilon$ do not affect lower and upper limits, as in the proof of \cref{prop:nef-degeneration-orders}. This gives $\underline a(\pi^*\theta)\geq\underline a(\theta)$ and $\overline a(\pi^*\theta)\geq\overline a(\theta)$. Finally, a proper holomorphic map of degree one between compact Riemann surfaces is an isomorphism, which settles the case $n=1$.
\end{proof}

\subsection{Strict changes under one blowup}\label{sec:nef-strict-blowups}

A blowup changes the amount of smooth volume near the vanishing set of a coefficient. A product equation turns that change into a strict increase of the degeneration order. Let $B=\mathbb P^m$, $m\geq2$, and normalize the Fubini--Study form $H$ by $\int_BH^m=1$. Let $\pi:\widetilde B\to B$ be the blowup of the affine origin. The incidence projection $\rho:\widetilde B\to\mathbb P^{m-1}$ gives $L=\rho^*\omega_{\rm FS}$; we also write $H$ for $\pi^*H$. The incidence embedding in $\mathbb P^m\times\mathbb P^{m-1}$ makes $\kappa=H+L$ K\"ahler.

\begin{proposition}\label{prop:nef-strict-blowup-orders}\index{Nef degeneration!strict blowup jump}
For every choice of integers $m\geq2$ and $p,k\geq1$, there is a single smooth-center blowup of a projective $(m+p)$-fold and a smooth semipositive nonclosed background whose common orders downstairs and after pullback are
\[
 a_{\rm down}=\max\{p-m/k,0\},\qquad a_{\rm up}=p-1/k.
\]
The orders exist on both sides. The increase can be any positive integer in varying dimension. The same family can lose positive limiting volume while both orders remain zero.
\end{proposition}

\begin{proof}
Put $W_0=\mathbb P^p$, with normalized Fubini--Study form $\omega$ satisfying $\int_{W_0}\omega^p=1$. Write $s=|z|^2$ on the affine chart of $B$. Choose a smooth $Q(s)\geq0$ equal to $s^k$ near zero, positive for $s>0$, and constant and positive for large $s$; it defines a smooth function $q$ on $B$. Put $n=m+p$, choose $b>0$, and define
\[
 \theta=bH+q\omega\quad\text{on }B\times W_0,\qquad
 \widetilde\theta=bH+(q\circ\pi)\omega
 \quad\text{on }\widetilde B\times W_0.
\]
The map $\Pi=\pi\times\mathrm{id}$ is the single blowup of the smooth center $\{0\}\times W_0$, and $\widetilde\theta=\Pi^*\theta$. Fix the K\"ahler directions $h=H+\omega$, $\widetilde h=\kappa+\omega$ and smooth positive probability densities $\mu=H^m\wedge\omega^p$, $\widetilde\mu=\kappa^m\wedge\omega^p/(2^m-1)$. Both backgrounds are smooth semipositive and nonclosed.

Write $\nu_{\rm down}=H^m$ and $\nu_{\rm up}=\kappa^m/(2^m-1)$ for the base probability forms, whose normalization is checked below. For either product write its perturbation as $\xi=A+g\omega$. Downstairs $A=(b+\epsilon)H$; upstairs $A=(b+\epsilon)H+\epsilon L$; in both cases $g=q+\epsilon$, with pullback understood upstairs. The form $A$ is K\"ahler on the relevant base. Use $\nu$ for the corresponding base form, and put $M=\int A^m$, $J=\int g^{-p}\nu$. The smooth equation \cref{thm:tw-smooth} and closed Stokes give $u$ with
\[
 \alpha=A+\ddc u>0,\qquad \alpha^m=(M/J)g^{-p}\nu.
\]
\Needspace{7\baselineskip}
The radial cumulative masses of the base forms are
\begin{equation}\label{eq:nef-radial-reference-masses}
 \begin{split}
 \Xi_{\rm down}(s)&=\left(\frac{s}{1+s}\right)^m,\\
 \Xi_{\rm up}(s)&=\frac{\left(1+\frac{s}{1+s}\right)^m-1}{2^m-1}.
 \end{split}
\end{equation}
To check these formulas and the constants $M$, write $t=\log|z|^2$ on the punctured affine chart. For a radial potential $F(t)$, direct polar integration of its Hessian gives
\[
 \int_{s_1<|z|^2<s_2}(\ddc F(\log|z|^2))^m
 =F'(\log s_2)^m-F'(\log s_1)^m.
\]
The potentials of $H$ and $L$ are $\log(1+\mathrm{e}^t)$ and $t$. Thus $xH+yL$ has endpoint slopes $y$ and $x+y$, and its total mass is $(x+y)^m-y^m$. The cases $(x,y)=(1,0),(1,1)$ give \eqref{eq:nef-radial-reference-masses}; the exceptional divisor and hyperplane at infinity have zero smooth volume. The form
\[
 \Omega=\alpha+g\omega=\xi+\ddc u>0,\qquad
\Omega^n=\binom{n}{m}\frac{M}{J}\nu\wedge\omega^p
\]
has constant volume density. The lower bound for all tests in \cref{thm:gaud-volume-dual-formula} is attained by
\[
 \gamma_*=g^{-p/(n-1)}\Omega.
\]
Indeed its $(n-1)$st power is
\begin{equation}\label{eq:nef-balanced-product-power}
 \gamma_*^{n-1}
 =\binom{n-1}{m}g^{-1}\alpha^m\wedge\omega^{p-1}
 +\binom{n-1}{m-1}\alpha^{m-1}\wedge\omega^p.
\end{equation}
The first coefficient multiplies a top-degree form on the base, and every factor in the second term is closed; hence both terms are closed. Thus $\gamma_*$ is balanced and Gauduchon. \Cref{thm:gaud-volume-dual-formula} gives the exact identities
\begin{equation}\label{eq:nef-strict-blowup-volumes}
 \mathcal V_\mu(\theta+\epsilon h)
 =\binom{n}{m}\frac{(b+\epsilon)^m}{J_{\rm down}},
 \qquad
 \mathcal V_{\widetilde\mu}(\widetilde\theta+\epsilon\widetilde h)
 =\binom{n}{m}
 \frac{(b+2\epsilon)^m-\epsilon^m}{J_{\rm up}},
\end{equation}
where $J_{\rm down}=\int_B(q+\epsilon)^{-p}\nu_{\rm down}$ and $J_{\rm up}=\int_{\widetilde B}(q\circ\pi+\epsilon)^{-p}\nu_{\rm up}$.

Near zero, differentiation of \eqref{eq:nef-radial-reference-masses} gives
\begin{equation}\label{eq:nef-radial-measure-asymptotics}
 \mathrm{d}\Xi_{\rm down}=(m s^{m-1}+O(s^m))\,\mathrm{d} s,\qquad
 \mathrm{d}\Xi_{\rm up}=\left(\frac{m}{2^m-1}+O(s)\right)\mathrm{d} s.
\end{equation}
\Needspace{10\baselineskip}
For a density $(a s^{d-1}+O(s^d))\,\mathrm{d} s$ and $q=s^k$, the singular integral has the following leading behavior as $\epsilon\downarrow0$:
\begin{equation}\label{eq:nef-radial-singular-integral}
 \begin{cases}
 \displaystyle\frac{a}{k}B(d/k,p-d/k)\epsilon^{d/k-p}
     &kp>d,\\[4pt]
 \displaystyle\frac{a}{k}\log(1/\epsilon)&kp=d,\\[4pt]
 \text{a finite positive limit}&kp<d.
 \end{cases}
\end{equation}
Here $B(u,v)=\int_0^\infty t^{u-1}(1+t)^{-u-v}\,\mathrm{d} t$. In the power case substitute $s=\epsilon^{1/k}v$, then $t=v^k$. At equality the integrand on $\epsilon^{1/k}\ll s\ll1$ is $a/s$. The $O(s^d)$ error is negligible: on a sufficiently small fixed interval its relative size is arbitrarily small, while the complementary integral is bounded. For $kp<d$, monotone convergence applies to the integrable limiting density. Since the numerators of \eqref{eq:nef-strict-blowup-volumes} tend to $\binom{n}{m}b^m>0$, formulas \eqref{eq:nef-radial-measure-asymptotics}--\eqref{eq:nef-radial-singular-integral}, with $d=m$ downstairs and $d=1$ upstairs, give the stated orders.

For $m=k=2$, $p=1$, the denominators satisfy $J_{\rm down}\sim\log(1/\epsilon)$ and $J_{\rm up}\sim(\pi/3)\epsilon^{-1/2}$; the volumes are respectively asymptotic to $3b^2/\log(1/\epsilon)$ and $(9b^2/\pi)\epsilon^{1/2}$. The orders change from $0$ to $1/2$. For $m=2$, $k=p=1$, the downstairs denominator has a finite positive limit, while $J_{\rm up}\sim(2/3)\log(1/\epsilon)$, so positive limiting volume downstairs becomes $9b^2/[2\log(1/\epsilon)]$ upstairs, with order zero on each side. For an arbitrary integer $j\geq1$, choose $m=j+2$, $p=j+1$, $k=1$. Then $p<m$ gives a positive downstairs limit, while the orders are $0$ and $j$.
\end{proof}

\subsection{Positive-current subtraction and divisorial residuals}\label{sec:nef-divisorial-residuals}

Another way to change the order of degeneration is to subtract a positive current from the background. If $\theta-\beta$ differs from a positive current by a $\ddc$-exact term, then every Gauduchon pairing of $\theta$ dominates the corresponding pairing of $\beta$. When the residual $\beta$ is still nef, its perturbed volumes are therefore smaller, and its orders are at least those of $\theta$. The main example is the subtraction of the curvature form of an effective divisor. Such a subtraction can preserve the order, raise it by an integer or by a fraction, or produce unequal lower and upper orders; and the residual need not be nef.

\Needspace{15\baselineskip}
\begin{proposition}\label{prop:nef-positive-current-subtraction}\index{Nef degeneration!positive-current subtraction}
Let $\theta,\beta$ be smooth nef real $(1,1)$-forms on a compact complex manifold $X$ of dimension $n$, and suppose that
\[
 S=(\theta-\beta)+\ddc\varphi\geq0
\]
is a positive $(1,1)$-current for some integrable function $\varphi$. For every Hermitian form $h$, every smooth strictly positive probability form $\mu$ and every $\epsilon>0$,
\begin{equation}\label{eq:nef-residual-volume-comparison}
 \mathcal V_\mu(\beta+\epsilon h)\leq \mathcal V_\mu(\theta+\epsilon h).
\end{equation}
Consequently $\underline a(\beta)\geq\underline a(\theta)$ and $\overline a(\beta)\geq\overline a(\theta)$.
\end{proposition}

\begin{proof}
For every Gauduchon metric $\gamma$, distributional integration by parts gives
\[
 \int_X(\theta-\beta)\wedge\gamma^{n-1}
 =\langle S,\gamma^{n-1}\rangle\geq0.
\]
Hence, for each test metric, the quotient for $\theta+\epsilon h$ is at least the quotient for $\beta+\epsilon h$, and taking infima gives $\lambda_\mu(\beta+\epsilon h)\leq\lambda_\mu(\theta+\epsilon h)$. Both root volumes are nonnegative, so raising to the $n$th power gives \eqref{eq:nef-residual-volume-comparison}. Dividing the logarithms by the negative number $\log\epsilon$ reverses the inequality and proves both comparisons of orders.
\end{proof}

For an effective real divisor $D$ with curvature form $\eta_D$, the Poincar\'e--Lelong formula gives $[D]=\eta_D+\ddc\varphi$ for a quasi-plurisubharmonic function $\varphi$. Thus \cref{prop:nef-positive-current-subtraction} applies to $\beta=\theta-\eta_D$, with $S=[D]$, whenever $\beta$ is nef. If $T=\theta+\ddc u$ is a positive current with $T-S\geq0$, then $T-S=\beta+\ddc(u-\varphi)$, so $\beta$ is the background of the positive current that remains after $S$ is removed from $T$.

\begin{proposition}\label{prop:nef-specified-residual-examples}
Subtracting the curvature form of an effective divisor from a smooth nef background, with a smooth nef residual, can preserve the common order, raise it by any integer between $1$ and the dimension or by a fraction, or produce a residual whose lower and upper orders differ. The residual can also fail to be nef.
\end{proposition}

\begin{proof}
\emph{Integer jumps.} On $E^n$, where $E$ is an elliptic curve, let $\omega_i$ be the pullback of a flat area form of integral one from the $i$th factor, and set $h=\sum_i\omega_i$ and $\mu=\omega_1\wedge\cdots\wedge\omega_n$. The divisor $F_i$ obtained by fixing a point in the $i$th factor has $\omega_i$ as a smooth curvature form. Indeed, on that elliptic curve any curvature form of the corresponding line bundle differs from $\omega_i$ by $\ddc$ of a solution of a zero-mean Poisson equation, so the metric can be adjusted on the curve and then pulled back. For $\theta=h$ and $D=\sum_{i=1}^kF_i$, the residual $\beta=\sum_{i=k+1}^n\omega_i$ is smooth and nef. The perturbed forms are closed, positive and of constant density with respect to $\mu$, so the zero potential solves their equations, and \cref{thm:gaud-volume-dual-formula} gives
\begin{equation}\label{eq:nef-integer-residual-jump}
 \mathcal V_\mu(\theta+\epsilon h)=n!(1+\epsilon)^n,\qquad
 \mathcal V_\mu(\beta+\epsilon h)=n!\epsilon^k(1+\epsilon)^{n-k}.
\end{equation}
These examples realize each integer jump $1,\ldots,n$. A nonzero divisor with jump zero is given by $\theta=2h$ and $D=F_1$, since then the residual is still strictly positive.

\emph{A fractional jump.} Use the background $\beta=b\omega_1+q\omega_2$ of \cref{prop:nef-diagonal-fractional-and-zero}, where $q=|\zeta|^{2m}$ near its unique zero and $m>1$. Choose $c>0$ and put $\theta=\beta+c\omega_2>0$ and $D=c(E\times\{w_0\})$, whose smooth curvature form is $c\omega_2$. The original background is strictly positive, so both its orders are zero, whereas the residual has common order $1-1/m$.

\emph{Unequal orders.} Let $\beta_{\mathrm{res}}$ be the oscillatory rank-one background of \cref{prop:nef-oscillatory-orders}, and put
\[
 \theta=\beta_{\mathrm{res}}+ch>0,\qquad
 D=c(\{z_0\}\times E)+c(E\times\{w_0\}),\qquad
 \eta_D=ch.
\]
The background $\theta$ has both orders zero: since $\theta\geq ch$, every test quotient is bounded below by a positive constant, while one fixed test gives a finite upper bound. Its smooth nef residual is exactly $\beta_{\mathrm{res}}$, whose lower and upper orders differ. In each of these examples, the residual plus $[D]$ is a positive current in the smooth $\ddc$-class of $\theta$, and its divisorial part is exactly $[D]$ because the residual is smooth. We do not need this current to have minimal singularities; that case is treated in \cref{cor:nef-boundary-minimal-removal}.

\emph{A residual that is not nef.} Blow up $\mathbb P^2$ at two distinct points. Let $H$ be the pullback of the class of the normalized Fubini--Study form, let $E_1,E_2$ be the exceptional curves, and let $L$ be the strict transform of the line through the two points. The pullback $\theta$ of the Fubini--Study form is smooth, semipositive and nef. Subtract curvature forms of $D=E_1+E_2$. The residual $\beta$ has class $H-E_1-E_2$, which is the class of $L$, and the Poincar\'e--Lelong formula for the total transform of the line shows that it contains the positive current $[L]$. The intersection numbers $H^2=1$, $H\cdot E_i=0$, $E_i^2=-1$ and $E_1\cdot E_2=0$ give
\[
 \int_L\beta=(H-E_1-E_2)^2=-1.
\]
If $\beta$ were nef, then, for a Hermitian form $h$, integrating smooth positive representatives $\beta+\epsilon h+\ddc u_\epsilon$ over $L$ would give $-1+\epsilon\int_Lh\geq0$ for all $\epsilon>0$, which is impossible.
\end{proof}

For a prime divisor $D$ and $\delta>0$, consider the infimum of the generic Lelong numbers along $D$ of the positive currents $\theta+\delta h+\ddc u$; its limit as $\delta\downarrow0$ is the minimal multiplicity of $\theta$ along $D$. For a nef background all minimal multiplicities vanish, since every $\theta+\delta h$ has a smooth positive representative. Removing the divisor defined by these multiplicities therefore leaves a nef background unchanged. Removing the divisorial part of a particular positive current in the class of $\theta$ is a different operation, even when that current has minimal singularities.

\begin{corollary}\label{cor:nef-boundary-minimal-removal}
On the ruled surface $X$ of \cref{lem:ruled-explicit-nef}, keep its nef curvature form $\theta$, its section $D$ at infinity and its K\"ahler form $\kappa=\theta+\pi^*\sigma_0$. For every smooth positive probability form $\mu$,
\[
 \mathcal V_\mu(\theta+\epsilon\kappa)=2\epsilon(1+\epsilon),
 \qquad \mathcal V_\mu(\epsilon\kappa)=2\epsilon^2.
\]
The only positive current of the form $\theta+\ddc u$ is $[D]$. It has minimal singularities, and removing its whole divisor leaves the zero background and changes the common order from one to two.
\end{corollary}
\begin{proof}
\Cref{lem:ruled-explicit-nef} gives the uniqueness of $[D]$ and the nefness of $\theta$, and the calculation in the proof of \cref{prop:nonrecovery-zero-residual-volume} gives $\int_X\theta^2=0$. Put $B=\pi^*\sigma_0$. Then $B^2=0$, and
\[
 \int_X\theta\wedge B=\int_D B=1,
\]
by the Poincar\'e--Lelong formula, Stokes' theorem and the fact that the section $D$ maps isomorphically onto the elliptic base. The closed nef formula following \cref{thm:gaud-volume-nef-limit} now computes both volumes by expanding the squares, for every $\mu$. Every $\theta$-plurisubharmonic function $u$ satisfies $\theta+\ddc u=[D]$, so any two such functions differ by a pluriharmonic function, hence by a constant; in particular each has minimal singularities. Since $\theta$ is itself a curvature form of $\mathcal O_X(D)$, removing the whole divisor leaves the zero background. The two volume polynomials give the orders one and two.
\end{proof}

\subsection{Integer orders in fixed projective cones}
\label{sec:nef-projective-integer-cones}
\index{Nef degeneration!nonclosed projective cones}

The variable-rank examples show that nefness alone does not force an integral order. Strictly positive coefficients on fixed projective generators give a different regime. A logarithmic divisor potential compares their full Gauduchon quotient with a closed form, even when the background is nonclosed and a pullback loses rank along an exceptional divisor.

\Needspace{14\baselineskip}
\begin{lemma}\label{lem:nef-logarithmic-rank-comparison}
Let $X$ be a compact K\"ahler $n$-fold, $\Omega$ a K\"ahler form, and $B$ a smooth closed semipositive form. Suppose a smooth nef background $\theta$ satisfies
\[
 \theta\leq CB,\qquad \theta+\ddc\varphi\geq cB
\]
as currents, for $c,C>0$ and a global real $\varphi\in L^1(X)$. For every Hermitian $h$ and smooth strictly positive probability $\mu$, there are $k,K>0$ such that, for $0<\epsilon\leq1$,
\begin{equation}\label{eq:nef-logarithmic-quotient-comparison}
 k\lambda_\mu(B+\epsilon\Omega)
 \leq\lambda_\mu(\theta+\epsilon h)
 \leq K\lambda_\mu(B+\epsilon\Omega).
\end{equation}
If $B$ has constant rank $R$, then for $L,U>0$,
\[
 L\epsilon^{n-R}\leq\mathcal V_\mu(\theta+\epsilon h)
                       \leq U\epsilon^{n-R}.
\]
In particular the common order is $n-R$.
\end{lemma}
\begin{proof}
Choose $c_h,C_h>0$ with $c_h\Omega\leq h\leq C_h\Omega$. For every smooth positive Gauduchon test $\gamma$, distributional Stokes gives
\[
 \int_X\ddc\varphi\wedge\gamma^{n-1}
       =\langle\varphi,\ddc\gamma^{n-1}\rangle=0.
\]
The lower current bound therefore gives the lower numerator inequality with $k=\min(c,c_h)$; the upper bound gives the upper one with $K=\max(C,C_h)$. Divide by the same strictly positive $D_\mu(\gamma)$ and take infima over all Gauduchon tests. This proves \eqref{eq:nef-logarithmic-quotient-comparison} without a rank hypothesis.

When $B$ has constant rank, its $R$ positive eigenvalues relative to $\Omega$ are uniformly bounded above and away from zero. With $A_\epsilon=B+\epsilon\Omega$, compactness gives
\begin{equation}\label{eq:nef-projective-pointwise-volume}
 l\epsilon^{n-R}\mu\leq A_\epsilon^n
                         \leq u\epsilon^{n-R}\mu
\end{equation}
for $l,u>0$. Matrix arithmetic--geometric mean gives, for every test,
\[
 A_\epsilon\wedge\gamma^{n-1}
 \geq\left(\frac{A_\epsilon^n}{\gamma^n}\right)^{1/n}\gamma^n.
\]
Thus $\lambda_\mu(A_\epsilon)\geq l^{1/n}\epsilon^{(n-R)/n}$. The positive closed form $\gamma=A_\epsilon$ is itself an admissible test. Writing $A_\epsilon^n=f_\epsilon\mu$ gives
\[
 \begin{aligned}
 D_\mu(A_\epsilon)&=\int_X f_\epsilon^{(n-1)/n}\mu\\
 &\geq l^{(n-1)/n}\epsilon^{(n-R)(n-1)/n},\\
 \int_X A_\epsilon^n&\leq u\epsilon^{n-R}.
 \end{aligned}
\]
Its quotient is at most $ul^{-(n-1)/n}\epsilon^{(n-R)/n}$. Raise both bounds to the $n$th power and apply the comparison. The calculation includes $R=0$ and $n=1$, where $D_\mu=1$. Fixed positive factors vanish after division of their logarithms by $\log\epsilon$.
\end{proof}

\subsubsection{Linear blowups and the fixed generators}
For a plain factor $Z_i=\mathbb P^{m_i}$, $m_i\geq1$, let $H_i$ be normalized Fubini--Study and set $\kappa_i=H_i$. A blown factor has the form
\[
 Z_i=\operatorname{Bl}_{\mathbb P^{\ell_i}}\mathbb P^{m_i},
 \qquad m_i\geq2,\quad 0\leq\ell_i\leq m_i-2,
 \qquad r_i=m_i-\ell_i-1.
\]
Suppress the index temporarily. In coordinates $[z_0:\cdots:z_\ell:w_0:\cdots:w_r]$, its incidence model is
\[
 Z=\{(x,[t])\in\mathbb P^m\times\mathbb P^r:
                         w_at_b=w_bt_a\text{ for all }a,b\}.
\]
Let $p:Z\to\mathbb P^m$ and $q:Z\to\mathbb P^r$ be the projections. In a chart $t_j=1$, the equations $w_a=t_aw_j$ are the charts of the blowup of $(w_0,\ldots,w_r)$. Writing $w=\lambda t$ identifies $q$ locally as a $\mathbb P^{\ell+1}$-bundle, with fiber coordinates $[z:\lambda]$. Thus $Z$ is smooth, $q$ is a submersion, and the incidence inclusion is an embedding. Set
\[
 H=p^*\omega_m,\qquad Q=q^*\omega_r,\qquad \kappa=H+Q.
\]
The form $\kappa$ is K\"ahler as the restricted product form, whereas $Q$ has constant rank $r$.

Let $E=\{w=0\}$ be the exceptional divisor. The coordinate section $w_j=\lambda t_j$ has divisor $E+q^{-1}\{t_j=0\}$, so $\mathcal O(E)=p^*\mathcal O(1)\otimes q^*\mathcal O(-1)$. Give it the product metric of curvature $H-Q$. For its canonical section, $\psi_E=\log|s_E|^2$ is locally integrable and the curvature identity of \cref{sec:divisorial-gaud-volume} gives
\begin{equation}\label{eq:nef-projective-exceptional-potential}
 \ddc\psi_E=[E]-H+Q,\qquad
 H+\ddc(\psi_E/2)=\tfrac12(\kappa+[E])\geq\kappa/2.
\end{equation}

\Needspace{9\baselineskip}
On $X=\prod_i Z_i$, put $n=\sum_i m_i$ and pull back all factor forms. Fix smooth strictly positive functions $f_i$ and, on blown factors, $g_i$ on $X$. These functions may depend on every product coordinate. The cone consists of
\begin{equation}\label{eq:nef-fixed-projective-cone}
 \theta=\sum_i\bigl((a_i+b_if_i)H_i+(c_i+d_ig_i)Q_i\bigr),
 \qquad a_i,b_i,c_i,d_i\geq0,
\end{equation}
with the $Q_i$ terms and their parameters omitted on plain factors. All coefficient functions and generators are fixed. Every member is smooth semipositive, hence nef for every Hermitian perturbation direction. Call $H_i$ active if $a_i+b_i>0$, and $Q_i$ active if $c_i+d_i>0$. Define
\begin{equation}\label{eq:nef-projective-effective-rank}
 R_i(\theta)=
 \begin{cases}
 m_i,&H_i\text{ active},\\
 r_i,&H_i\text{ inactive and }Q_i\text{ active},\\
 0,&\text{both inactive},
 \end{cases}
 \qquad R(\theta)=\sum_i R_i(\theta).
\end{equation}
The middle case occurs only on blown factors.

\Needspace{9\baselineskip}
\begin{theorem}\label{thm:nef-projective-integer-orders}
Every fixed-cone member in \eqref{eq:nef-fixed-projective-cone} has common order $n-R(\theta)$. For arbitrary smooth positive $\mu$ and $h$, its variational volume has two-sided bounds for $0<\epsilon\leq1$:
\begin{equation}\label{eq:nef-projective-power-bounds}
 L\epsilon^{n-R(\theta)}
 \leq\mathcal V_\mu(\theta+\epsilon h)
 \leq U\epsilon^{n-R(\theta)},
\end{equation}
with $L,U>0$ depending on the fixed background and data.
\end{theorem}
\begin{proof}
Let $B_\theta$ be the sum of $\kappa_i$ on $H_i$-active blocks and $Q_i$ on $Q_i$-only blocks, with zero on inactive blocks. It is closed semipositive of constant rank $R(\theta)$. Bounded coefficients and $H_i,Q_i\leq\kappa_i$ give $\theta\leq CB_\theta$.

For every active blown $H_i$, let $u_i=\min_X(a_i+b_if_i)>0$ and put
\[
 \varphi=\sum_{\substack{i\text{ blown}\\H_i\text{ active}}}
                          (u_i/2)\psi_{E_i}.
\]
By \eqref{eq:nef-projective-exceptional-potential}, $u_iH_i$ becomes at least $(u_i/2)\kappa_i$, with a nonnegative divisor current. The remaining $(a_i+b_if_i-u_i)H_i$ and all $Q_i$ terms are nonnegative. On plain active blocks and $Q_i$-only blocks, use the positive minimum of their active coefficient. Taking the minimum over the active blocks gives $\theta+\ddc\varphi\geq cB_\theta$. For the zero member take $\varphi=0$ and any $c>0$. Apply \cref{lem:nef-logarithmic-rank-comparison} with $\Omega=\sum_i\kappa_i$.
\end{proof}

Nonclosed members occur whenever a fixed generator derivative is nonzero, since
\[
 \mathrm d\theta=\sum_i b_i\,\mathrm df_i\wedge H_i
                  +\sum_{i\text{ blown}}d_i\,\mathrm dg_i\wedge Q_i.
\]
Activating that generator alone prevents cancellation. On a blown factor, a small positive bump perturbation of one varying in a $q$-fiber has $\mathrm dg\wedge Q\ne0$, since $r\geq1$. On $\mathbb P^p\times\mathbb P^q$, $p,q\geq1$, a positive $f$ varying on the second factor gives nonclosed $fH_1$ of order $q$. Adding a smooth $\ddc$ potential preserves the exterior derivative.

\subsubsection{Modification equality}
\Needspace{11\baselineskip}
\begin{proposition}\label{prop:nef-projective-modification-equality}
Every product modification obtained by replacing plain factors with the indicated linear blowups preserves the order of a fixed-cone member. More generally, for every modification $\pi:Y\to X$ with $Y$ a compact K\"ahler manifold, the pullback has common order $n-R(\theta)$ and two-sided power bounds with that exponent, for arbitrary smooth positive probability and Hermitian perturbation on $Y$.
\end{proposition}
\begin{proof}
For product linear blowups, pull back the old coefficient functions. The new $Q_i$ coefficients are zero. Active plain $H_i$ become active blown $H_i$ of the same effective rank $m_i$, and inactive blocks stay inactive. Apply \cref{thm:nef-projective-integer-orders} on both products.

For general K\"ahler $Y$, pull back the explicit nonnegative smooth remainders and divisor-section identities above. Each nonzero divisor section remains nonzero under the surjective modification. Its pulled-back squared logarithm is locally integrable and gives the curvature of its effective zero divisor, so
\[
 \pi^*\theta\leq C\pi^*B_\theta,\qquad
 \pi^*\theta+\ddc(\varphi\circ\pi)\geq c\pi^*B_\theta.
\]
Use the rank-free part of \cref{lem:nef-logarithmic-rank-comparison} with a K\"ahler form $\Omega_Y$. The pulled-back background is semipositive and hence nef. Put $R=R(\theta)$. Pullback cannot increase rank, so $(\pi^*B_\theta)^{R+1}=0$, while
\[
 \int_Y(\pi^*B_\theta)^R\wedge\Omega_Y^{n-R}>0.
\]
For $R>0$ this is positive on the nonempty isomorphism locus; for $R=0$ it is the K\"ahler volume. The closed-background formula following \cref{thm:gaud-volume-nef-limit} gives, for the chosen probability $\mu_Y$,
\[
 \mathcal V_{\mu_Y}(\pi^*B_\theta+\epsilon\Omega_Y)
 =\sum_{k=0}^{R}\binom nk\epsilon^{n-k}
       \int_Y(\pi^*B_\theta)^k\wedge\Omega_Y^{n-k}.
\]
Every coefficient is nonnegative and the coefficient at $k=R$ is positive. Thus this polynomial has two-sided bounds by $\epsilon^{n-R}$ for $0<\epsilon\leq1$. The quotient comparison gives the same bounds for $\pi^*\theta+\epsilon h_Y$. This also proves equality with the downstairs order. No pullback of a positive probability is assumed to remain strictly positive on the exceptional locus.
\end{proof}

\Needspace{13\baselineskip}
\subsubsection{Specified effective divisors}
Fix the pullback Fubini--Study metrics on $\mathcal O(H_i)$ and $\mathcal O(Q_i)$. Allow tensor products of their integer powers, including negative powers. For finite nonnegative real sums of zero divisors of nonzero holomorphic sections of these bundles, write
\[
 \begin{gathered}
 D=\sum_\nu t_\nu\operatorname{div}(s_\nu),\quad t_\nu\geq0,\\
 \eta_D=\sum_i(\delta_iH_i+\rho_iQ_i),\qquad
 \Phi_D=\sum_\nu t_\nu\log|s_\nu|^2.
 \end{gathered}
\]
The metrics fix the representatives, and $[D]=\eta_D+\ddc\Phi_D$. Omit $Q_i$ terms on plain factors.

\Needspace{8\baselineskip}
\begin{lemma}\label{lem:nef-projective-effective-divisor-range}
The exact coefficient range of this real divisor family is
\[
 \delta_i\geq0\text{ on plain factors};\qquad
 \delta_i\geq0,\quad\delta_i+\rho_i\geq0
                         \text{ on blown factors}.
\]
\end{lemma}
\begin{proof}
For one nonzero holomorphic section $s$, choose a point where it is nonzero and outside all exceptional divisors, then fix the other factor coordinates. A projective line through that point in a $q_i$-fiber has $H_i$-degree one and $Q_i$-degree zero. The restriction of $s$ is nonzero, so its $H_i$ exponent is nonnegative. On a plain factor the line argument gives the same inequality.

On a blown factor, write the chosen point as $([z:w],[w])$ with $w\ne0$, and choose a linear map $A$ with $Aw=z$. The section of $q_i$ given by $[t]\mapsto([At:t],[t])$ contains the point. Both $H_i$ and $Q_i$ have degree one on any of its projective lines. Such a line exists because $r_i\geq1$, and the original holomorphic section $s$ restricts nontrivially to it. Its degree gives the nonnegative sum of the two exponents. Nonnegative real sums preserve these inequalities.

Conversely, let $F_i$ be a $Q_i$-hyperplane divisor. A permitted pair is represented by $\delta_iE_i+(\delta_i+\rho_i)F_i$, since $E_i$ has curvature $H_i-Q_i$. Plain coefficients are represented by hyperplanes. These are allowed effective real sums and give every coefficient tuple in the stated range.
\end{proof}

\Needspace{13\baselineskip}
\begin{theorem}\label{thm:nef-projective-residual-equality}
Let $\theta$ have the fixed-cone parameters above, and let $D$ be in the specified divisor family. Assume
\[
 a_i\geq\delta_i\text{ for all }i,\qquad
 c_i\geq\rho_i\text{ for every blown }i.
\]
Then $\beta=\theta-\eta_D$ is in the same cone, and its common order satisfies
\[
 a(\beta)-a(\theta)=\sum_i\bigl(R_i(\theta)-R_i(\beta)\bigr)\geq0.
\]
Equality holds if and only if
\begin{align*}
 a_i-\delta_i+b_i&>0&&\text{whenever }a_i+b_i>0,\\
 c_i-\rho_i+d_i&>0&&\text{whenever }a_i=b_i=0,\quad c_i+d_i>0.
\end{align*}
Here $a$ denotes the common order; the second condition concerns blown factors only.
\end{theorem}
\begin{proof}
The residual parameters $(a_i-\delta_i,b_i,c_i-\rho_i,d_i)$ are nonnegative, so $\beta$ is smooth semipositive and nef. Also
\[
 (\theta-\beta)+\ddc\Phi_D=[D]\geq0,
\]
which is the actual positive-current comparison of \cref{prop:nef-positive-current-subtraction}. Apply \cref{thm:nef-projective-integer-orders} to both members. If an active $H_i$ disappears, its rank drops from $m_i$ to either $r_i<m_i$ or zero. If $H_i$ was absent, then $a_i=b_i=0$; the divisor range and residual conditions force $\delta_i=0$ and $\rho_i\geq0$. Such a block cannot gain $Q_i$. Therefore each rank loss is nonnegative and the losses cannot cancel. Their sum vanishes precisely when every initially leading generator remains active, giving the two conditions.
\end{proof}

The parameter inequalities ensure that this displayed representation stays in the fixed cone. They do not characterize every nef residual or every alternative representation by its generators.

\Needspace{10\baselineskip}
\begin{example}\label{ex:nef-nonclosed-projective-order-jumps}
On $Z=\operatorname{Bl}_{\mathbb P^\ell}\mathbb P^m$, choose $g>0$ with $\mathrm dg\wedge Q\ne0$. For $t>0$, the genuinely nonclosed positive background $\theta_t=tH+gQ$ has order zero. Removing $sE$, with its fixed curvature $s(H-Q)$ and $0\leq s\leq t$, gives
\[
 \beta_s=(t-s)H+(g+s)Q,\qquad
 \theta_t+\ddc(s\psi_E)=s[E]+\beta_s.
\]
Every residual is nef and nonclosed. Equality of orders holds exactly for $s<t$; at $s=t$ the jump is $\ell+1$. These examples realize each integer jump $1,\ldots,m-1$.
\end{example}
\begin{proof}
The lower bound $\theta_t\geq\min(t,\min g)\kappa>0$ verifies positivity. For $s<t$, $H$ is active and the effective rank is $m$; at $s=t$ only $Q$ is active, of rank $r=m-\ell-1$. The order formula gives zero and $m-r=\ell+1$, respectively. The exterior derivative of every residual is $\mathrm dg\wedge Q\ne0$. Vary $\ell$ from zero to $m-2$.
\end{proof}

At the endpoint the chosen current contains the actual divisor $t[E]$ before removal. The perturbative minimal divisorial multiplicities are nevertheless zero, since all residual backgrounds admit smooth positive perturbations. Strict positivity of the fixed coefficient functions, the specified cone-preserving divisor inequalities and a K\"ahler source for the general modification statement are essential hypotheses of this construction. The general problem in \cref{qu:nef-degeneration-exponent} remains open.

\subsection{A maximum-rank comparison criterion}
\label{sec:nef-maximum-rank-comparison}
\index{Nef degeneration!maximum-rank comparison}

The constant-rank hypothesis in \cref{lem:nef-logarithmic-rank-comparison} can be removed. For a closed semipositive form $B$, the intersection polynomial $\int_X(B+\epsilon\Omega)^n$ has a leading term of order $\epsilon^{n-R}$, where $R$ is the largest rank that $B$ attains anywhere on $X$. Combined with the comparison of quotients, this gives the order $n-R$ for every background that is comparable to $B$. The criterion applies to backgrounds that vanish on an open set, and to positive variable multiples of arbitrary closed semipositive generators.

\Needspace{12\baselineskip}
\begin{proposition}\label{prop:nef-maximum-rank-comparison}
Let $X$ be a compact K\"ahler manifold of dimension $n$, let $B$ be a smooth closed semipositive $(1,1)$-form, and put $R=\max_X\operatorname{rank}B$. Suppose a smooth nef background $\theta$ satisfies
\[
 \theta\leq CB,\qquad \theta+\ddc\varphi\geq cB
\]
as currents, for constants $c,C>0$ and a real function $\varphi\in L^1(X)$. For every smooth positive probability form $\mu$ and Hermitian form $h$, there are constants $L,U>0$ such that
\[
 L\epsilon^{n-R}\leq\mathcal V_\mu(\theta+\epsilon h)
                       \leq U\epsilon^{n-R},\qquad 0<\epsilon\leq1.
\]
Thus $\theta$ has common order $n-R$, without any constant-rank hypothesis.
\end{proposition}
\begin{proof}
Fix a K\"ahler form $\Omega$. The first part of \cref{lem:nef-logarithmic-rank-comparison}, which does not use the rank, shows that $\lambda_\mu(\theta+\epsilon h)$ and $\lambda_\mu(B+\epsilon\Omega)$ differ by bounded positive factors. Since $B^{R+1}=0$, the closed nef formula following \cref{thm:gaud-volume-nef-limit} gives
\[
 \mathcal V_\mu(B+\epsilon\Omega)
 =\sum_{k=0}^{R}\binom nk\epsilon^{n-k}
                         \int_X B^k\wedge\Omega^{n-k}.
\]
All coefficients are nonnegative. If $R>0$, then at a point where $B$ has rank $R$, some $R\times R$ minor of its coefficient matrix is positive and stays positive nearby, so $B^R\wedge\Omega^{n-R}$ has positive integral. If $R=0$, this integral is the K\"ahler volume. Hence the coefficient of $\epsilon^{n-R}$ is positive. For $0<\epsilon\leq1$, every other power of $\epsilon$ that occurs is at most $\epsilon^{n-R}$, so the polynomial is bounded below by its leading term and above by the sum of its coefficients times $\epsilon^{n-R}$. Raising the comparison of quotients to the $n$th power gives the asserted bounds, and taking logarithms gives the order.
\end{proof}

\Needspace{12\baselineskip}
\begin{corollary}\label{cor:nef-variable-semipositive-generators}
On a compact K\"ahler manifold $X$ of dimension $n$, fix finitely many smooth closed semipositive $(1,1)$-forms $B_j$ and smooth strictly positive functions $f_j$. For $t_j\geq0$, put
\[
 \theta=\sum_jt_jf_jB_j,\qquad
 B=\sum_{j:t_j>0}B_j,\qquad R=\max_X\operatorname{rank}B.
\]
The background $\theta$ is nef, has common order $n-R$, and satisfies the two-sided power bounds of \cref{prop:nef-maximum-rank-comparison}. For every modification $\pi:Y\to X$ with $Y$ a compact K\"ahler manifold, the pullback $\pi^*\theta$ has the same order and power bounds with the same exponent, for every smooth positive probability form and Hermitian perturbation on $Y$.
\end{corollary}
\begin{proof}
Semipositivity gives nefness. If some $t_j$ is positive, compactness and strict positivity of the functions $f_j$ give
\[
 \begin{gathered}
 cB\leq\theta\leq CB,\qquad
 c=\min_{j:t_j>0}\bigl(t_j\min_X f_j\bigr)>0,\\
 C=\max_{j:t_j>0}\bigl(t_j\max_X f_j\bigr),
 \end{gathered}
\]
and \cref{prop:nef-maximum-rank-comparison} applies with $\varphi=0$. If all parameters vanish, then $B=\theta=0$, $R=0$, any $c,C>0$ may be used, and the order is $n$.

These inequalities between smooth forms are preserved by pullback, and the rank of $\pi^*B$ is at most $R$. If $R>0$, the set where $B$ has rank $R$ is nonempty and open, so it meets the dense open set where $\pi$ is biholomorphic; hence the maximal rank of $\pi^*B$ is $R$. This is immediate if $R=0$. Since $\dim Y=n$, the same proposition applies on $Y$, with any choice of data there. No singular potential needs to be pulled back.
\end{proof}

The constants depend on the parameters $t_j$ and need not stay bounded as a positive parameter tends to zero. The rank of a sum records how the generators overlap, and it need not equal the sum of their ranks. Members of the family are not closed when $\sum_jt_j\,\mathrm df_j\wedge B_j\ne0$.

\Needspace{10\baselineskip}
\begin{example}\label{ex:nef-open-zero-rank-background}
Let $X=\mathbb P^1\times\mathbb P^q$, $q\geq1$, with projections $p_1,p_2$. Choose a smooth nonnegative function $b\not\equiv0$ on $\mathbb P^1$ that vanishes on a nonempty open set, and a smooth strictly positive nonconstant function $f$ on $\mathbb P^q$. With $\omega_1$ the normalized Fubini--Study form on $\mathbb P^1$, set
\[
 B=p_1^*(b\omega_1),\qquad \theta=(f\circ p_2)B.
\]
The smooth nef background $\theta$ is not closed, vanishes on a nonempty open set, and has common order $q$, with two-sided bounds by multiples of $\epsilon^q$.
\end{example}
\begin{proof}
The form $b\omega_1$ is closed, because the first factor is a curve, and semipositive. Its pullback has rank one where $b>0$ and rank zero where $b=0$, so its maximal rank is one. Since $n=q+1$, \cref{cor:nef-variable-semipositive-generators} gives the order $q$. A nonconstant smooth function $f$ on the connected second factor has $\mathrm df\ne0$ somewhere. At a point of the product where $b>0$ and $\mathrm df\neq0$, the horizontal two-form $B$ wedges nontrivially with the vertical one-form $p_2^*\mathrm df$. Thus $\mathrm d\theta=p_2^*\mathrm df\wedge B\ne0$.
\end{proof}

These are geometric sufficient conditions for an integral common order, in the direction of \cref{qu:nef-degeneration-exponent}. They do not cover backgrounds without a closed semipositive comparison form, coefficient functions that are allowed to vanish, or the comparison of orders under residuals or non-K\"ahler modifications in general.

\subsection{Flat anticanonical curves}\label{subsec:nef-flat-anticanonical-curves}

A flat neighborhood of a divisor gives a closed semipositive form whose curvature is confined to a thin annulus around the divisor. On a surface this form has rank at most one. The maximum-rank criterion then applies to its positive variable multiples, even when the divisor line has no moving positive multiple. We first construct the flat neighborhood for a square-zero anticanonical curve.

\begin{lemma}\label{lem:flat-elliptic-derivative-inverse}
Let $D$ be an elliptic curve with a flat Hermitian metric, let $\vartheta$ be a parallel unit $(1,0)$-form, and let $v$ be its dual vector field. For every unitary flat line bundle $L$ on $D$ and every smooth $L$-valued $(0,1)$-form $u$, there is a unique smooth section $f$ orthogonal to the parallel sections such that
\[
 \bar\partial_L f=\iota_v\partial_Lu.
\]
If $\Pi_L^{0,1}$ is the orthogonal projection onto the parallel $L$-valued $(0,1)$-forms, then, for every flat Sobolev norm $H^k$,
\[
 \|f\|_{H^k}=\|(I-\Pi_L^{0,1})u\|_{H^k}
 \leq\|u\|_{H^k}.
\]
The bound is independent of the holonomy of $L$.
\end{lemma}
\begin{proof}
Put $P=\iota_v\partial_L$. Flatness gives $P=\nabla_v$, $P^*=-\nabla_{\bar v}$, and $\Delta_L=P^*P=PP^*$ on functions and $(0,1)$-forms. Let $G_q$ be the inverse of $\Delta_L$ on the orthogonal complement of its kernel in degree $q$, extended by zero on that kernel. The operators $P$ and $P^*$ commute with $G_q$. Set
\[
 f=\bar\partial_L^*G_1Pu.
\]
The form $Pu$ is orthogonal to the parallel forms, so $\bar\partial_Lf=Pu$; also $f$ is orthogonal to the parallel sections. Uniqueness follows from the kernel of $\bar\partial_L$. Moreover,
\[
 \|f\|_{L^2}^2
 =\langle Pu,G_1Pu\rangle
 =\langle u,(I-\Pi_L^{0,1})u\rangle.
\]
Covariant differentiation along the global parallel real frame commutes with these operators. Applying the identity to every derivative proves the $H^k$ assertion. In particular, this argument never divides by a lower bound for a nonzero eigenvalue; the derivative in $Pu$ cancels that small eigenvalue.
\end{proof}

For the neighborhood argument, let $S$ be a complex surface, possibly noncompact, and let $D\subset S$ be a smooth compact connected anticanonical divisor with $D^2=0$. Adjunction makes $D$ elliptic. Its normal line $N$ has degree zero, hence admits a unitary flat metric. Let $\pi:N\to D$ be its projection, $N^\times$ its punctured total space, and $J_N$ its complex structure. The reciprocal anticanonical section gives a meromorphic two-form $\Omega_S$, nowhere zero on $S\setminus D$, with simple pole along $D$. Denote its nonzero residue by $\vartheta_D$.

Choose parallel unitary local frames of $N$ with fiber coordinate $w$. On $N^\times$ put
\[
 \xi=w\partial_w,\qquad \Theta=\frac{\mathrm dw}{w},\qquad
 \vartheta=\pi^*\vartheta_D,\qquad
 z=-\log|w|^2,\qquad \Omega_0=\Theta\wedge\vartheta.
\]
These objects are global. Give $D$ the flat metric making $\vartheta_D$ unit, and let $v$ be the parallel vector field dual to $\vartheta_D$. Equip $N^\times$ with the flat cylindrical metric $g$ determined by the parallel coframe $\vartheta,\Theta$, and write $\nabla$ for its connection. The end $z\to+\infty$ approaches the zero section.

\begin{lemma}\label{lem:anticanonical-cylindrical-identification}
There is a smooth tubular identification $\Upsilon$ from a neighborhood of the zero section in $N$ to a neighborhood of $D$ in $S$, fixing $D$, such that each fiber disc is holomorphic for $\Upsilon^*J_S$ and
\[
 |\nabla^k(\Upsilon^*\Omega_S-\Omega_0)|_g
 =O(\mathrm e^{-z/2})\qquad(k\geq0).
\]
\end{lemma}
\begin{proof}
Start with a normal exponential map for a smooth Hermitian metric near $D$. Its differential is complex linear along $D$ and induces the identity on the normal quotient. In a parallel local frame and a coordinate $z_0$ with $\vartheta_D=\mathrm dz_0$, choose holomorphic surface coordinates $Z,W$ with $D=\{W=0\}$. The pullbacks under this first map satisfy
\[
 Z=z_0+b(z_0,\bar z_0)w+O(|w|^2),\qquad
 W=w+O(|w|^2).
\]
The residue normalization writes $\Omega_S=(1+Wa(Z,W))\mathrm dW/W\wedge\mathrm dZ$ with $a$ holomorphic. Substitution gives $\Omega_S-\Omega_0=O(|w|)$ in the logarithmic coframe. The possible term containing $\Theta\wedge\mathrm dw$ vanishes. Every derivative by $\partial_{z_0},\partial_{\bar z_0},w\partial_w,\bar w\partial_{\bar w}$ preserves this bound, proving the stated cylindrical estimate for the first map.

It remains to make the fiber discs holomorphic without changing its first jet. On a fiber disc $\Delta_{p,\rho}\subset N_p$, vary the map by $\Phi_Y(q)=\exp_q(Y_q)$. Identify its target tangent spaces with a fixed complex vector space over $p$ by Chern parallel transport along the radial paths and the short exponential paths. In these identifications its antiholomorphic differential is
\[
 \mathcal F_{p,\rho}(Y)
 =\tfrac12\bigl(\mathrm d\Phi_Y+J_S\mathrm d\Phi_YJ_N\bigr)
       |_{T\Delta_{p,\rho}}.
\]
For a fixed Hölder exponent $0<\nu<1$, use rescaled $C^{k,\nu}$ norms on the discs. The Cauchy--Green operator $T_\rho$ gives the normalized right inverse
\[
 C_\rho a=T_\rho a-T_\rho a(0)
                  -w\partial_wT_\rho a(0).
\]
Restrict its domain to $(0,1)$-forms vanishing at the center. Its image consists of fields with zero value and first jet there; on this image $\bar\partial$ and $C_\rho$ are inverse. The differential of $\mathcal F_{p,\rho}$ at zero has the form $\bar\partial+K_{p,\rho}$, with
\[
 \|K_{p,\rho}C_\rho\|\leq C\rho,\qquad
 \|\mathcal F_{p,\rho}(0)\|\leq C\rho^2.
\]
Indeed, the complex structures agree along the zero section, while pulling a one-form back by $w\mapsto\rho w$ supplies the additional factor $\rho$. The inverse is $C_\rho(I+K_{p,\rho}C_\rho)^{-1}$ for small $\rho$. The quantitative implicit-function theorem therefore gives a correction with zero first jet and $Y=O(|w|^2)$.

Compactness of $D$ makes the choice uniform in $p$. The normalized Cauchy operator is equivariant under the unitary changes of fiber frame; the intrinsic transports and uniqueness of the small solution make the corrections agree on overlaps. The parametric implicit-function theorem gives smooth dependence on $p$, and interior elliptic regularity on slightly smaller fiber discs gives smoothness in the fiber. The fixed first jet then gives the quadratic ordinary error, with its cylindrical derivative bounds. Composing the first map with this correction preserves the estimate and proves the lemma.
\end{proof}

\begin{theorem}\label{thm:flat-anticanonical-neighborhood}
\index{Anticanonical bundle!flat elliptic neighborhood}\index{Divisor!flat neighborhood}
Let $S$ be a complex surface and $D\subset S$ a smooth compact connected anticanonical divisor with $D^2=0$. The line bundle $\mathcal O_S(D)$ is unitary flat on a neighborhood of $D$.
\end{theorem}
\begin{proof}
Use the cylindrical identification of \cref{lem:anticanonical-cylindrical-identification}, and omit its pullback notation. A form is horizontal if contraction with $\xi$ and $\bar\xi$ vanishes. For a horizontal form define
\[
 \mathrm d_H\eta=\mathrm d\eta
       -\Theta\wedge\mathcal L_\xi\eta
       -\bar\Theta\wedge\mathcal L_{\bar\xi}\eta.
\]
The degree-$m$ horizontal lift of a form $\eta_m$ valued in $N^{-m}$ is $w^m\eta_m$ in a parallel frame. Direct differentiation identifies $\mathrm d_H$ with its flat covariant derivative. Since the fibers are holomorphic,
\[
 \iota_{\bar\xi}\Omega_S=0,\qquad
 \Omega_S=\Theta\wedge\alpha+\beta,
 \qquad\alpha=\iota_\xi\Omega_S,
\]
with $\alpha,\beta$ horizontal. Closedness gives $\mathcal L_{\bar\xi}\alpha=\mathcal L_{\bar\xi}\beta=0$ and $\mathcal L_\xi\beta=\mathrm d_H\alpha$.

Fix $0<\delta<1/2$, and let $U_T=\{z\geq T\}$. For tensors define the weighted norm
\[
 \|\eta\|_{k,\delta,T}^2
 =\sum_{j=0}^k\int_{U_T}\mathrm e^{2\delta(z-T)}
                    |\nabla^j\eta|_g^2\,\mathrm dV_g.
\]
Let $\mathcal H^{k+1}_{\delta,T}$ be the functions of finite such norm with $\mathcal L_{\bar\xi}f=0$, and let $\mathcal D^k_{\delta,T}$ be the complex two-forms of finite norm with $\mathrm d\gamma=0$ and $\iota_{\bar\xi}\gamma=0$. We construct an operator
\[
 Q_T:\mathcal D^k_{\delta,T}\longrightarrow\mathcal H^{k+1}_{\delta,T},
 \qquad \|Q_T\gamma\|_{k+1,\delta,T}\leq C_{k,\delta}\|\gamma\|_{k,\delta,T},
\]
with constant independent of $T$, satisfying
\begin{equation}\label{eq:anticanonical-horizontal-inverse}
 \mathrm d_H(Q_T\gamma\,\vartheta)
 =\gamma-\Theta\wedge\iota_\xi\gamma,
 \qquad Q_T(\mathrm df\wedge\vartheta)=f.
\end{equation}

Write $\gamma=\Theta\wedge\alpha+\beta$. Its defining equations give $\mathcal L_{\bar\xi}\gamma=0$ and, in the fiberwise homogeneous expansion, $m\beta_m=\mathrm d_H\alpha_m$. Only $m\geq1$ occur: the finite positive weighted norm excludes the zero and negative Laurent modes. Decompose
\[
 \alpha_m=-u_m\vartheta+v_m,
 \qquad u_m\in\Omega^0(D,N^{-m}),\quad
 v_m\in\Omega^{0,1}(D,N^{-m}).
\]
Apply \cref{lem:flat-elliptic-derivative-inverse} to get $w_m$ with $\bar\partial w_m=\iota_v\partial v_m$ and $\|w_m\|_{H^k}\leq\|v_m\|_{H^k}$. Put $f_m=-(u_m+w_m)/m$. Then
\[
 \mathrm d_H(f_m\vartheta)=\beta_m,\qquad
 \|mf_m\|_{H^k}\leq C\|\alpha_m\|_{H^k}.
\]
Also $\|\bar\partial f_m\|_{H^k}\leq C\|\beta_m\|_{H^k}$. Flatness gives the same bound for $\partial f_m$ by integration by parts, and hence controls the pure horizontal derivatives of order $k+1$. A vertical derivative of a homogeneous mode multiplies it by $m$, so every derivative involving a vertical direction is controlled by the preceding bound for $mf_m$. Circle orthogonality and integration of the weight give the required estimate for $f=\sum_{m\geq1}f_m$: the squared radial factor is $\mathrm e^{-mz}$ and $m-2\delta>0$.

For Sobolev coefficients, regularize each mode by the flat base heat operator, which commutes with the defining differential equations. The uniform estimates and orthogonality make the finite sums Cauchy in the target space. Passing to their limit proves the first identity in \eqref{eq:anticanonical-horizontal-inverse}. If $\gamma=\mathrm df\wedge\vartheta$, then $\alpha_m=mf_m\vartheta$, so its $v_m$ and normalized $w_m$ vanish and the construction returns $f_m$. This proves the second identity.

Put $\eta=\Omega_S-\Omega_0$. Choose an integer $k\geq4$. The cylindrical estimate bounds its weighted $H^{k+2}$ norm on $U_{T-1}$ and its first $k+2$ uniform derivatives by a number $a_T\to0$. For small $f\in\mathcal H^{k+1}_{\delta,T}$, the map $\Phi_f(\zeta)=\mathrm e^{f(\zeta)}\zeta$ is a $C^2$ diffeomorphism onto its image and sends $U_T$ into $U_{T-1}$. Its restriction to every fiber is holomorphic. Thus $\Phi_f^*\Omega_S-\Omega_0$ belongs to the actual domain of $Q_T$. Since
\[
 \Phi_f^*\Omega_0-\Omega_0=\mathrm df\wedge\vartheta,
\]
the desired equation becomes
\[
 Q_T(\Phi_f^*\Omega_S-\Omega_0)=f+E_T(f)=0,
 \qquad E_T(f)=Q_T(\Phi_f^*\eta).
\]
On a fixed sufficiently small ball, weighted Sobolev multiplication and composition on the flat cylinder give
\[
 \|E_T(f)\|_{k+1,\delta,T}\leq Ca_T,\qquad
 \|E_T(f)-E_T(g)\|_{k+1,\delta,T}
 \leq Ca_T\|f-g\|_{k+1,\delta,T}.
\]
The constants are independent of $T$: translation in the cylindrical direction is an isometry, the weight is normalized at $T$, and the shifts of $z$ are uniformly bounded. Differentiating the pullback along the segment between $f$ and $g$ takes one Lie derivative of the fixed form $\eta$ and at most one derivative of $f-g$; the gain of one derivative in $Q_T$ restores the indicated norm. For $T$ sufficiently large, $-E_T$ is a contraction of the closed ball of radius $2Ca_T$ into itself. Let $f$ be its fixed point.

The first identity in \eqref{eq:anticanonical-horizontal-inverse} makes the horizontal component of $\Phi_f^*\Omega_S$ vanish. Consequently $\Theta\wedge\Phi_f^*\Omega_S=0$. The pulled-back volume form is nowhere zero and of type $(2,0)$, so $\Theta$ is of type $(1,0)$ for $\Phi_f^*J_S$. Since $\mathrm dz=-(\Theta+\bar\Theta)$ and $\mathrm d\Theta=0$, the function $z$ is pluriharmonic for that complex structure.

Return to $S$ through $A=\Upsilon\circ\Phi_f$. On each fiber, boundedness extends $f$ holomorphically across zero. Its positive-mode expansion gives $f(0)=0$. The uniform small bound on $f$ lets Rouch\'e's theorem compare $w\mathrm e^{f(w)}-a$ with $w-a$ on a fixed smaller circle, for every sufficiently small $a$. Thus the image of $A$ contains a smaller punctured neighborhood of $D$. The function $r=z\circ A^{-1}$ is $C^2$ and pluriharmonic there, hence smooth by elliptic regularity. If $w_S$ is a holomorphic defining function for $D$, the fixed first normal jet of $\Upsilon$ and the bounded factor $\mathrm e^f$ give
\[
 C^{-1}|w|\leq|w_S\circ A|\leq C|w|.
\]
Thus $r+\log|w_S|^2$ is bounded and pluriharmonic off $D$. Both signs have plurisubharmonic extensions across the divisor; their sum is zero almost everywhere and hence everywhere. The extension is therefore pluriharmonic. It follows that $F=-r$, with its logarithmic pole on $D$, satisfies $\ddc F=[D]$ near $D$. Locally, the norm assigning $\exp(F)$ to the squared norm of the defining section is a smooth flat metric, since $F-\log|w_S|^2$ is smooth pluriharmonic. These norms agree on overlaps and give the asserted unitary flat metric.
\end{proof}

\begin{proposition}\label{prop:flat-divisor-rank-one-curvature}
Let $D$ be a smooth nonempty compact divisor in a complex surface $S$, and suppose $\mathcal O_S(D)$ is unitary flat near $D$. It has a smooth semipositive metric on $S$ whose curvature $B$ has rank at most one and vanishes near $D$ and outside a relatively compact neighborhood of $D$. If $S$ is compact K\"ahler, then $\max_S\operatorname{rank}B=1$.
\end{proposition}
\begin{proof}
Let $s_D$ be the defining section and $h_0$ a flat metric on a neighborhood $U$ of $D$. Put $F=\log|s_D|_{h_0}^2$. Choose a relatively compact tubular neighborhood $U_0\Subset U$. Since $F$ tends uniformly to $-\infty$ along $D$ and is finite on $\partial U_0$, adding a constant makes $F>0$ near $\partial U_0$. Choose a smooth concave function $\chi$ with $\chi(t)=t+c$ for $t\leq-1$ and $\chi(t)=0$ for $t\geq0$. Define $\widetilde F=\chi(F)$ on $U_0$ and extend it by zero away from $U_0$. On $S\setminus D$ prescribe
\[
 |s_D|_h^2=\mathrm e^{\widetilde F}.
\]
Near $D$ this gives $h=\mathrm e^c h_0$, so the frame metric extends smoothly and positively across $D$. It also matches smoothly where the cutoff is constant. Poincar\'e--Lelong gives $B=[D]-\ddc\widetilde F$, and off $D$ one has
\[
 B=-\chi''(F)\,\mathrm dF\wedge\dc F\geq0.
\]
This smooth closed form has rank at most one and the asserted vanishing regions. On a compact K\"ahler surface with K\"ahler form $\kappa$, its mixed integral is
\[
 \int_SB\wedge\kappa=\int_D\kappa>0.
\]
Thus it is not identically zero and its maximum rank is one.
\end{proof}

\Needspace{15\baselineskip}
\begin{corollary}\label{cor:flat-divisor-variational-order}
\index{Nef degeneration!flat divisor curvature}
Let $S$ be a compact K\"ahler surface and $D$ a smooth nonempty divisor whose line bundle is unitary flat near $D$. Let $B$ be the curvature from \cref{prop:flat-divisor-rank-one-curvature}. For every fixed smooth $f>0$, every Hermitian form $h$ and every smooth positive probability $\mu$, there are $L,U>0$ such that
\[
 L\epsilon\leq\mathcal V_\mu(fB+\epsilon h)
              \leq U\epsilon,\qquad0<\epsilon\leq1.
\]
For a K\"ahler perturbation $\kappa$, the closed background itself satisfies
\[
 \mathcal V_\mu(B+\epsilon\kappa)
 =2\epsilon\int_D\kappa+\epsilon^2\int_S\kappa^2.
\]
There are choices of $f$ for which $fB$ is nonclosed.
\end{corollary}
\begin{proof}
The inequalities $(\min f)B\leq fB\leq(\max f)B$ and \cref{prop:nef-maximum-rank-comparison} give the two-sided bounds. The polynomial follows from the closed-volume formula and $B^2=0$. At a point where $B$ has rank one, choose a smooth positive $f$ whose differential has a nonzero component tangential to the kernel of $B$. Then $\mathrm d(fB)=\mathrm df\wedge B\ne0$ there.
\end{proof}

\Needspace{12\baselineskip}
\begin{proposition}\label{prop:nine-point-anticanonical-curvature}
\index{Anticanonical bundle!nine-point blowup}
Let $S$ be obtained from $\mathbb P^2$ by nine successive point blowups, each centered on the strict transform of a fixed smooth cubic. Its final strict transform $D$ is a smooth elliptic anticanonical curve with $D^2=0$. The anticanonical bundle has a smooth semipositive metric with rank-one curvature as above. It is semiample exactly when $N=\mathcal O_D(D)$ is torsion. If $N$ has exact order $m$, then $|mD|$ is basepoint-free and no smaller positive multiple moves. If $N$ is non-torsion, then $|kD|=\{kD\}$ for every $k\geq1$.
\end{proposition}
\begin{proof}
A blowup at a smooth point of the cubic replaces its class by its pullback minus the exceptional curve, exactly as for the anticanonical class. The strict transform remains isomorphic to the cubic, and its self-intersection decreases by one. After nine steps $D^2=9-9=0$. The flat neighborhood theorem and curvature construction apply.

For $k\geq1$, use the restriction sequences
\[
 0\longrightarrow\mathcal O_S((k-1)D)
 \longrightarrow\mathcal O_S(kD)
 \longrightarrow N^k\longrightarrow0.
\]
A nontrivial degree-zero line on an elliptic curve has no section: any nonzero section would have a zero divisor of degree zero and trivialize the line. If $N$ is non-torsion, induction gives $h^0(S,\mathcal O_S(kD))=h^0(S,\mathcal O_S)=1$, proving the singleton systems and failure of semiampleness.

If $N$ has exact order $m$, then $H^0(D,N^k)=H^1(D,N^k)=0$ for $1\leq k<m$, the latter by Serre duality on the elliptic curve. Since $S$ is rational, $H^1(S,\mathcal O_S)=0$. The same sequences give $h^0(S,kD)=1$ and $H^1(S,\mathcal O_S(kD))=0$ for $0\leq k<m$. The constant section of $N^m\simeq\mathcal O_D$ therefore lifts to a section of $\mathcal O_S(mD)$ nonzero everywhere on $D$. Together with $s_D^m$, it generates that line globally. Thus $|mD|$ is basepoint-free, and no smaller multiple moves.

Non-torsion configurations exist with nine distinct centers. Fix the first eight points on the cubic and vary the ninth. Under the unchanged identification with that elliptic curve, the normal line is $\mathcal O_D(3)\otimes\mathcal O_D(-p_1-\cdots-p_9)$. Varying $p_9$ parametrizes its degree-zero Picard group by the elliptic Abel map. The torsion subgroup is countable, while the curve is uncountable; excluding the first eight points still leaves non-torsion choices. Hence \cref{cor:flat-divisor-variational-order} supplies nonclosed backgrounds of order one even when the associated anticanonical line is not semiample.
\end{proof}

\subsection{Flat divisors and Hartogs extension}
\label{subsec:flat-divisor-complements}

All the nine-point blowups of \cref{prop:nine-point-anticanonical-curvature} carry backgrounds of order one, so the order of degeneration does not distinguish the torsion case, in which a multiple of $D$ moves in a basepoint-free pencil, from the non-torsion case, in which no multiple of $D$ moves. The difference is visible in the complex geometry near $D$. On the tube where the curvature form of \cref{prop:flat-divisor-rank-one-curvature} vanishes, the logarithm of the flat norm of the defining section is pluriharmonic. Its regular level sets are therefore Levi-flat, that is, real hypersurfaces foliated by complex hypersurfaces. These complex leaves are governed by the unitary holonomy of the normal bundle: finite holonomy gives compact fibers of a local holomorphic function, whereas infinite holonomy makes every leaf dense and forces every holomorphic function on the punctured tube to be constant. In both cases the complement of $D$ is not Stein.

In this subsection $S$ is a complex manifold of dimension $n\geq2$ and $D\subset S$ is a smooth compact connected divisor. Suppose that $L=\mathcal O_S(D)$ has a unitary flat metric $h$ on a neighborhood of $D$. Write $s_D$ for its defining section and $N=L|_D$ for the normal bundle, with the induced metric and flat connection, and let $\chi\colon\pi_1(D)\to\mathrm U(1)$ be the holonomy character of $N$, which describes parallel transport around closed loops. In parallel unit frames, $s_D$ is represented by holomorphic functions $W_j$ satisfying $W_j=t_{jk}W_k$ with locally constant $t_{jk}\in\mathrm U(1)$. The flat connection defines on each circle bundle $\{v\in N:|v|=r\}$ a foliation whose leaves are given locally, in parallel frames, by $\{w=\text{constant}\}$; we call it the flat horizontal foliation.

\Needspace{16\baselineskip}
\begin{lemma}\label{lem:flat-divisor-foliated-tube}
There are $R>0$, a neighborhood $U$ of $D$ and a smooth diffeomorphism
\[
 \Upsilon\colon\{v\in N:|v|<R\}\longrightarrow U
\]
that is the identity on the zero section and preserves the defining function: in each parallel normal frame, $W_j(\Upsilon(p,w))=w$. In particular, $|s_D|_h\circ\Upsilon=|w|$. For $0<r<R$, the set
\[
 M_r=\{x\in U:|s_D(x)|_h=r\}
\]
is a compact connected Levi-flat hypersurface, and $\Upsilon$ maps the leaves of the flat horizontal foliation of $\{|w|=r\}$ onto the complex leaves of $M_r$.
\end{lemma}
\begin{proof}
Choose a smooth normal exponential map $\Upsilon_0$ whose first normal derivative is complex linear; the derivative of $s_D$ identifies the normal bundle with $L|_D$, and we use this identification to fix that first derivative. Parallel transport along the normal rays identifies $\Upsilon_0^*L$ with the pullback of $N$. In compatible parallel frames one then has
\[
 W_j(\Upsilon_0(p,a))=a+O(|a|^2).
\]
For each $p$, this is a smooth map between real two-dimensional disks whose derivative at zero is the identity. The inverse function theorem with parameters, uniformly over the compact base $D$, gives a unique small $a=a(p,w)$ with $W_j(\Upsilon_0(p,a))=w$. It depends smoothly on $(p,w)$, and its first derivative in $w$ at $w=0$ is the identity. On an overlap of frames, $a$ and $w$ change by the same constant unitary factor, so uniqueness makes these corrections agree. After decreasing $R$, the map $(p,w)\mapsto(p,a(p,w))$ is a diffeomorphism of a tube onto its image, and its composition with $\Upsilon_0$ is the required map $\Upsilon$.

Off $D$, the function $F=\log|s_D|_h^2$ is pluriharmonic, and its radial derivative in the tube is nonzero. Hence every $M_r$ is a smooth hypersurface; it is compact and connected, being a circle bundle over $D$. Locally its complex tangent space is $\ker\mathrm dW_j$, so it is Levi-flat, with complex leaves the hypersurfaces $\{W_j=\text{constant}\}$. The identity $W_j\circ\Upsilon=w$ pulls this distribution back to $\{\mathrm dw=0\}$ in parallel frames, which is the flat horizontal foliation of the circle bundle. The map $\Upsilon$ is only a smooth tubular identification; its base projection need not be holomorphic.
\end{proof}

\begin{proposition}\label{prop:flat-divisor-finite-holonomy}
In this setting, $N$ is torsion as a holomorphic line bundle if and only if its unitary holonomy has finite image. If $N$ has exact order $m$, then a sufficiently small tube $U$ admits a proper surjective holomorphic map
\[
 \phi\colon U\longrightarrow\{z\in\mathbb C:|z|<R^m\}
\]
with connected fibers and $\operatorname{div}(\phi)=mD$. Its fibers over nonzero values are smooth compact complex hypersurfaces.
\end{proposition}
\begin{proof}
If $\chi^m=1$, a parallel nowhere-zero section trivializes $N^m$ holomorphically. Conversely, if a holomorphic section $e$ trivializes $N^m$, then $\log|e|^2$ is pluriharmonic on the compact connected manifold $D$, hence constant by the maximum principle. In parallel holomorphic frames the coefficient of $e$ is a holomorphic function of constant modulus, hence constant. Therefore $e$ is parallel and $\chi^m=1$. Consequently the least such $m$ is both the order of $N$ as a torsion line bundle and the order of the finite holonomy group.

By \cref{lem:flat-divisor-foliated-tube}, the tube $U$ retracts onto $D$. Hence the flat bundle $L^m$ on $U$ has trivial holonomy and a global parallel section $e_U$ of unit length. Define $\phi=s_D^m/e_U$. Then $\phi$ is holomorphic, $\operatorname{div}(\phi)=mD$ and $|\phi|=|s_D|_h^m$. The inverse image of a compact subset of the target lies in a smaller closed normal disk bundle and is closed there, which proves properness, and the tubular model shows that every value in the disk is attained.

Choose local parallel unit frames $e_j$ of $L$ with $e_j^m=e_U$, so that $\phi=W_j^m$. In the tubular model, the fiber over a nonzero value consists, over each point of $D$, of the $m$ roots of that value. Parallel transport acts on these roots through $\chi$, whose image has exact order $m$ and therefore acts transitively. Thus this fiber is a connected covering of $D$, as a real manifold. Since $\mathrm d\phi=mW_j^{m-1}\mathrm dW_j$, it is a smooth complex hypersurface, and it is compact because $\phi$ is proper. The fiber over zero has connected support $D$, which completes the proof that all fibers are connected.
\end{proof}

\begin{theorem}\label{thm:flat-divisor-holomorphic-rigidity}
\index{Hartogs extension!non-torsion flat divisor}
\index{Levi-flat hypersurface!flat divisor tube}
Suppose $N$ is not torsion. Then every complex leaf of $M_r$ is dense in $M_r$, and every holomorphic function on $U\setminus D$ is constant:
\[
 H^0(U\setminus D,\mathcal O_S)=\mathbb C.
\]
Suppose moreover that $S$ is connected. For every compact set $K\subset S\setminus D$ such that $(S\setminus D)\setminus K$ is connected, every holomorphic function on $(S\setminus D)\setminus K$ is constant, and in particular extends holomorphically to $S\setminus D$. Similarly, for every connected open neighborhood $V$ of $D$, every holomorphic function on $V\setminus D$ is constant and extends to $S\setminus D$.
\end{theorem}
\begin{proof}
By \cref{prop:flat-divisor-finite-holonomy}, $\chi(\pi_1(D))$ is an infinite subgroup of $\mathrm U(1)$. Such a subgroup is dense, since a proper closed subgroup of the circle is finite, as one sees by lifting it to a closed subgroup of $\mathbb R$ containing $2\pi\mathbb Z$. The circle bundle with its flat horizontal foliation is the suspension of $\chi$, that is, the quotient of $\widetilde D\times S^1$ by the diagonal action of $\pi_1(D)$, acting on $S^1$ through $\chi$. A leaf projects onto $D$, and its intersection with each fiber contains the holonomy orbit of any one of its points, which is dense in the fiber. In a local trivialization, horizontal transport maps the orbits in nearby fibers to each other by a fixed rotation, so the leaf is dense in the whole circle bundle. By \cref{lem:flat-divisor-foliated-tube}, every complex leaf is dense in $M_r$.

Let $f$ be holomorphic on $U\setminus D$. Its modulus attains a maximum on the compact set $M_r$. On the complex leaf through a maximum point, $|f|$ is bounded by this maximum and attains it, so the maximum modulus principle makes $f$ constant on the connected leaf. By density and continuity, $f$ equals the same constant $c$ on $M_r$. Near any point of $M_r$, choose holomorphic coordinates $(z,W_j)$ with $z\in\mathbb C^{n-1}$. For each fixed $z$, the holomorphic function $W_j\mapsto f(z,W_j)-c$ vanishes on an arc of the circle $|W_j|=r$, so it vanishes on the coordinate neighborhood by the one-variable identity theorem. Analytic continuation on the connected punctured tube gives $f\equiv c$.

A compact set $K\subset S\setminus D$ is disjoint from a sufficiently small tube around the compact divisor $D$. A holomorphic function on $(S\setminus D)\setminus K$ is therefore constant on this punctured tube, and by the identity theorem on the connected set $(S\setminus D)\setminus K$ it is constant everywhere; the constant function is the required extension. Similarly, every open neighborhood $V$ of $D$ contains a full tube, and $V\setminus D$ is connected, because paths in the connected manifold $V$ can be perturbed off the real codimension-two submanifold $D$. The identity theorem gives the assertion for $V$. The same path argument shows that $S\setminus D$ is connected.
\end{proof}

Recall that a complex manifold $Y$ is \emph{holomorphically separable} if its holomorphic functions separate points, and \emph{holomorphically convex} if the holomorphic hull
\[
 \widehat K_Y=\{y\in Y:|f(y)|\leq\max_K|f|\text{ for every }f\in\mathcal O(Y)\}
\]
of every compact set $K\subset Y$ is compact. A Stein manifold has both properties.

\begin{corollary}\label{cor:flat-divisor-nonstein}
\index{Stein manifold!flat divisor obstruction}
For any holonomy of $N$, the punctured tube $U\setminus D$ admits no smooth strictly plurisubharmonic function, and $S\setminus D$ is not Stein. In the non-torsion nine-point blowups of \cref{prop:nine-point-anticanonical-curvature}, every holomorphic function on the complement of $D$ is constant, and the Hartogs extension of \cref{thm:flat-divisor-holomorphic-rigidity} across compact holes holds.
\end{corollary}
\begin{proof}
The restriction of a smooth strictly plurisubharmonic function on $U\setminus D$ to a complex leaf of $M_r$ is strictly plurisubharmonic. At a maximum point on the compact set $M_r$, this restriction has a local maximum on the leaf, which has positive dimension because $n\geq2$; this contradicts the maximum principle.

In the finite-holonomy case, each fiber of the map of \cref{prop:flat-divisor-finite-holonomy} over a nonzero value is a connected compact complex hypersurface of positive dimension. By the maximum modulus principle, every holomorphic function on $S\setminus D$ is constant on it, so two distinct points of the fiber cannot be separated. In the infinite-holonomy case, \cref{thm:flat-divisor-holomorphic-rigidity} makes every holomorphic function constant on the punctured tube, which again contradicts holomorphic separability. In both cases $S\setminus D$ is not Stein. Finally, \cref{thm:flat-anticanonical-neighborhood} gives the flat neighborhood for the nine-point blowups, whose normal bundle is non-torsion by assumption, and \cref{thm:flat-divisor-holomorphic-rigidity}, also with $K=\varnothing$, gives the last assertion.
\end{proof}

The nine-point blowups show that a background with a well-defined order of degeneration can have very different holomorphic geometry near its null locus; \cref{app:cotangent-torsors} gives a second instance of this contrast, for K\"ahler classes on a ruled surface over an elliptic curve.


\begin{problem}\label{qu:nef-degeneration-exponent}
Which geometric conditions on a nonclosed nef background force a common degeneration order, or an integral one? When is equality attained in the modification and residual order comparisons?
\end{problem}
